\def\micro{\mu}
\documentclass[10pt,a4paper,oneside]{article}
\usepackage[left=2cm,right=2cm,top=2cm,bottom=2cm]{geometry}

../../global_color_scheme.tex
\usepackage[utf8]{inputenc}
\usepackage[english]{babel}
\usepackage{forloop}
\usepackage{amsmath}
\usepackage{amsfonts}
\usepackage{amssymb}
\usepackage{gensymb}
\usepackage{mdframed}
\usepackage{graphicx}
\usepackage{tikz}
\usetikzlibrary{arrows,automata,shapes}
\usepackage[siunitx]{circuitikz}

\usepackage[colorlinks=true,linkcolor=blue,urlcolor=black,bookmarksopen=true]{hyperref}
\usepackage{bookmark}
\usepackage{hyperref}
\usepackage{sepfootnotes}

\usepackage[digital,srcmeas]{circdia}

\usepackage{float}
\floatstyle{boxed} 
\restylefloat{figure}

\title{Libre Silicon process specification}
\date{\today} 
\author{David Lanzendörfer}
\makeindex

\newcounter{ct}
\def\CrossSectionOnly{0.3}
\def\CrossAndTopSection{0.2}
\def\CrossAndTopSectionBig{0.3}
\def\VLSILayout{0.4}

\DeclareMathOperator\erfc{erfc}

\setlength{\parindent}{0pt} % get rid of annoying indents

\begin{document}
\maketitle

\begin{abstract}
	Copyright © 2017 LANCEVILLE TECHNOLOGY GROUP CO., LIMITED. All rights reserved. \\

	This process is licensed under the Libre Silicon public license; you can redistribute it and/or modify it under the terms of the Libre Silicon public license
	as published by the Libre Silicon alliance, either version 1 of the License, or (at your option) any later version.

	This design is distributed in the hope that it will be useful, but WITHOUT ANY WARRANTY; without even the implied warranty of MERCHANTABILITY or FITNESS FOR A PARTICULAR PURPOSE.
	See the Libre Silicon Public License for more details. \\

	This is the specification of the free silicon manufacturing standard for manufacturing the LibreSilicon standard logic cells\footnote{\url{https://github.com/chipforge/StdCellLib}} and related free technology nodes from the LibreSilicon project.
\end{abstract}

\newpage
\tableofcontents
\newpage
\section{CMOS in a nutshell}
This basic initial project is dedicated to the CMOS Technology only and for this reason two types of metal–oxide–semiconductor field-effect transistors (MOSFET) are required.

Historicaly, the first chips with MOSFETs on the mass market were p-channel MOSFETs in enhancement-mode.

\begin{center}
	\begin{figure}[h]
		\begin{center}
			\begin{circuitdiagram}{20}{20}
				\power{15}{18.5}{U}{}  % power above pmos
				\wire{15}{18}{15}{16}   % wire above pmos
				\trans{penh}{13}{14}{R}{}{} % pmos -> right
				\Voltarrow{14}{18}{10}{16}{u}{$-V_{GS}$}
				\wire{15}{12}{15}{8}   % wire below pmos
				\resis{15}{5}{V}{$R_D$}{}  % resistor on drain
				\wire{15}{1}{15}{2}   % wire below pmos
				\ground{15}{0.5}{D}  % ground below resistor
				\othersrc[\sigsym{-rec}]{o}{5.5}{15.5}{H}{}{signal}
				\pin{9.5}{15.5}{R}{}	% pin in
				\ground{2.5}{0.5}{D}  % ground below signal source
				\wire{2.5}{1}{2.5}{15.5}   % wire below signal
				\junct{15}{10}   % dot
				\wire{15}{10}{16}{10}   % wire before out
				\pin{17}{10}{R}{out}	% pin out
			\end{circuitdiagram}
		\end{center}
		\caption{enhancement-mode PMOS transistor use-case}
	\end{figure}
\end{center}

The sectional view of a PMOS transistor in silicon is being shown below
\begin{center}
	\begin{figure}[h]
		\begin{center}
			\begin{tikzpicture}[node distance = 3cm, auto, thick,scale=0.5, every node/.style={transform shape}]
				% substrate
				\fill[YellowOrange] (0,0) rectangle (10,2);
				\node at (2,0.5) {Si (p-type)};

				% n-well
				\fill[Goldenrod] (1.25,0.75) rectangle (8.25,2);
				\node at (5.75,1) {N-Well};

				% body
				\fill[ProcessBlue] (1.5,1) rectangle (3,2);
				\node at (2,1.5) {n+};
				% source
				\fill[RedOrange] (3.5,1) rectangle (5,2);
				\node at (4,1.5) {p+};
				% drain
				\fill[RedOrange] (6.5,1) rectangle (8,2);
				\node at (7,1.5) {p+};
				%% gate:
				% gate oxide
				\fill[LightGray] (4.8,2) rectangle (6.7,2.1);
				% gate poly
				\fill[BrickRed] (4.8,2.1) rectangle (6.7,2.2);

				% metals ground
				\fill[DarkDarkGray] (2,2) rectangle (2.5,3);
				\fill[DarkDarkGray] (4,2) rectangle (4.5,3);
				\fill[DarkDarkGray] (1,3) rectangle (4.5,3.2); % connection pad GND

				\fill[DarkDarkGray] (5.5,2.2) rectangle (6,3);
				\fill[DarkDarkGray] (5.3,3) rectangle (6.2,3.2); % connection pad VG

				\fill[DarkDarkGray] (7,2) rectangle (7.5,3);
				\fill[DarkDarkGray] (6.8,3) rectangle (7.7,3.2); % connection pad VDD

				% isolation oxides:
				\fill[NormalGray] (1,2) rectangle (2,3);
				\fill[NormalGray] (2.5,2) rectangle (4,3);
				\fill[NormalGray] (4.5,2) rectangle (4.8,3);
				\fill[NormalGray] (4.8,2.2) rectangle (5.5,3);
				\fill[NormalGray] (6,2.2) rectangle (6.7,3);
				\fill[NormalGray] (6.7,2) rectangle (7,3);
				\fill[NormalGray] (7.5,2) rectangle (8.5,3);

				\node at (1.5,3.5) {Source};
				\node at (5.5,3.4) {Gate};
				\node at (7.5,3.4) {Drain};

				%field oxides:
				\fill[DarkGray] (0,2) rectangle (1,4);
				\fill[DarkGray] (8.5,2) rectangle (10,4);
				\fill[RedOrange] (0,1.5) rectangle (1,2);
				\node at (0.5,1.75) {p+};
				\fill[RedOrange] (8.5,1.5) rectangle (10,2);
				\node at (9.5,1.75) {p+};
			\end{tikzpicture}
		\end{center}
		\caption{Sectional view of a PMOS transistor}
	\end{figure}
\end{center}

Historicaly later, faster chips with MOSFETs on the mass market were marked as n-channel MOSFETs in enhancement mode also.

\begin{center}
	\begin{figure}[h]
		\begin{center}
			\begin{circuitdiagram}{20}{20}
				\power{15}{18.5}{U}{}  % power above resistor
				\wire{15}{17}{15}{18}   % wire above resistor
				\resis{15}{14}{V}{$R_D$}{}  % resistor on drain
				\wire{15}{11}{15}{8}   % wire between resistor and nmos
				\trans{nenh}{13}{6}{R}{}{} % nmos -> right
				\Voltarrow{10}{4}{14}{1}{d}{$+V_{GS}$}
				\wire{15}{1}{15}{4}   % wire below nmos
				\ground{15}{0.5}{D}  % ground below nmos
				\othersrc[\sigsym{rec}]{o}{5.5}{4.5}{H}{}{signal}
				\pin{9.5}{4.5}{R}{}	% pin in
				\ground{2.5}{0.5}{D}  % ground below signal source
				\wire{2.5}{1}{2.5}{4.5}   % wire below signal
				\junct{15}{10}   % dot
				\wire{15}{10}{16}{10}   % wire before out
				\pin{17}{10}{R}{out}	% pin out
			\end{circuitdiagram}
		\end{center}
		\caption{enhancement-mode NMOS transistor use-case}
	\end{figure}
\end{center}

The sectional view of a NMOS transistor in silicon is being shown here also.
\begin{center}
	\begin{figure}[h]
		\begin{center}
			\begin{tikzpicture}[node distance = 3cm, auto, thick,scale=0.5, every node/.style={transform shape}]
				% substrate
				\fill[YellowOrange] (0,0) rectangle (10,2);
				\node at (2,0.5) {Si (p-type)};
				% body
				\fill[RedOrange] (1.5,1) rectangle (3,2);
				\node at (2,1.5) {p+};
				% source
				\fill[ProcessBlue] (3.5,1) rectangle (5,2);
				\node at (4,1.5) {n+};
				% drain
				\fill[ProcessBlue] (6.5,1) rectangle (8,2);
				\node at (7,1.5) {n+};
				%% gate:
				% gate oxide
				\fill[LightGray] (4.8,2) rectangle (6.7,2.1);
				% gate poly
				\fill[BrickRed] (4.8,2.1) rectangle (6.7,2.2);
				% metals ground
				\fill[DarkDarkGray] (2,2) rectangle (2.5,3);
				\fill[DarkDarkGray] (4,2) rectangle (4.5,3);
				\fill[DarkDarkGray] (1,3) rectangle (4.5,3.2); % connection pad GND

				\fill[DarkDarkGray] (5.5,2.2) rectangle (6,3);
				\fill[DarkDarkGray] (5.3,3) rectangle (6.2,3.2); % connection pad VG

				\fill[DarkDarkGray] (7,2) rectangle (7.5,3);
				\fill[DarkDarkGray] (6.8,3) rectangle (7.7,3.2); % connection pad VDD

				% isolation oxides:
				\fill[NormalGray] (1,2) rectangle (2,3);
				\fill[NormalGray] (2.5,2) rectangle (4,3);
				\fill[NormalGray] (4.5,2) rectangle (4.8,3);
				\fill[NormalGray] (4.8,2.2) rectangle (5.5,3);
				\fill[NormalGray] (6,2.2) rectangle (6.7,3);
				\fill[NormalGray] (6.7,2) rectangle (7,3);
				\fill[NormalGray] (7.5,2) rectangle (8.5,3);

				%field oxides:
				\fill[DarkGray] (0,2) rectangle (1,4);
				\fill[DarkGray] (8.5,2) rectangle (10,4);
				\fill[RedOrange] (0,1.5) rectangle (1,2);
				\node at (0.5,1.75) {p+};
				\fill[RedOrange] (8.5,1.5) rectangle (10,2);
				\node at (9.5,1.75) {p+};

				\node at (1.5,3.5) {Source};
				\node at (5.5,3.4) {Gate};
				\node at (7.5,3.4) {Drain};
			\end{tikzpicture}
		\end{center}
		\caption{Sectional view of a NMOS transistor}
	\end{figure}
\end{center}

Both technologies, the older NMOS as the newer PMOS, have the same disadvantage. Every time, the transistor is switched on, the current between Drain and Source of the transistor is limited by the Resistor on Drain only. Higher currents here meaning also higher power consumption for the chip where the transistors are integrated also. If the transistors are switched off, now currents flows between Drain and Source anymore, the power consumption of the chip also goes low.
Et violà, the US-Patent with Number 3356858\footnote{https://www.google.com/patents/US3356858} changed the world and combines both technologies to the new complementary metal-oxide-semiconductor (CMOS) technology. Instead of every transistor is working against a weak resistor, the transistor works against a complementary switched-off transistor. With the Eyes of our antecessor CMOS doubles the transistor count, but contemporary chips all are build in CMOS.

\begin{center}
	\begin{figure}[h]
		\begin{center}
			\begin{circuitdiagram}{20}{20}
				\power{15}{18.5}{U}{}  % power above pmos 
				\wire{15}{16}{15}{18}   % wire above pmos
				\trans{penh}{13}{14}{R}{}{} % pmos -> right
				\Voltarrow{14}{18}{10}{16}{u}{$-V_{GS}$}
				\wire{15}{8}{15}{12}   % wire between pmos and nmos
				\trans{nenh}{13}{6}{R}{}{} % nmos -> right
				\Voltarrow{10}{4}{14}{1}{d}{$+V_{GS}$}
				\wire{15}{1}{15}{4}   % wire below nmos
				\ground{15}{0.5}{D}  % ground below nmos
				\othersrc[\sigsym{rec}]{o}{5}{10}{H}{}{signal}
				\pin{9}{10}{R}{}	% pin in
				\wire{9.5}{10}{10}{10}   % wire before gates
				\wire{10}{15.5}{10}{4.5}   % wire between gates
				\junct{10}{10}   % dot
				\ground{2}{0.5}{D}  % ground below signal source
				\wire{2}{1}{2}{10}   % wire below signal
				\junct{15}{10}   % dot
				\wire{15}{10}{16}{10}   % wire before out
				\pin{17}{10}{R}{out}	% pin out
			\end{circuitdiagram}
		\end{center}
		\caption{complementary PMOS and NMOS transistor couple use-case}
	\end{figure}
\end{center}

Below the sectional view of the inverter circuitry can be seen.
For the run through of this process we will use this cross section diagram as reference.
\begin{center}
	\begin{figure}[h]
		\begin{center}
			\begin{tikzpicture}[node distance = 3cm, auto, thick,scale=0.5, every node/.style={transform shape}]
				% substrate
				\fill[YellowOrange] (0,0) rectangle (20,2);
				\node at (2,0.5) {Si (p-type)};
				% n-well
				\fill[Goldenrod] (1.25,0.75) rectangle (8.25,2);
				\node at (5.75,1) {N-Well};
				% body
				\fill[ProcessBlue] (1.5,1) rectangle (3,2);
				\node at (2,1.5) {n+};
				% source
				\fill[RedOrange] (3.5,1) rectangle (5,2);
				\node at (4,1.5) {p+};
				% drain
				\fill[RedOrange] (6.5,1) rectangle (8,2);
				\node at (7,1.5) {p+};
				%% gate:
				% gate oxide
				\fill[LightGray] (4.8,2) rectangle (6.7,2.1);
				% gate poly
				\fill[BrickRed] (4.8,2.1) rectangle (6.7,2.2);
				% isolation oxides:
				\fill[NormalGray] (1,2) rectangle (2,3);
				\fill[NormalGray] (2.5,2) rectangle (4,3);
				\fill[NormalGray] (4.5,2) rectangle (4.8,3);
				\fill[NormalGray] (4.8,2.2) rectangle (5.5,3);
				\fill[NormalGray] (6,2.2) rectangle (6.7,3);
				\fill[NormalGray] (6.7,2) rectangle (7,3);
				\fill[NormalGray] (7.5,2) rectangle (8.5,3);

				%field oxides:
				\fill[DarkGray] (0,2) rectangle (1,4);
				\fill[DarkGray] (8.5,2) rectangle (11.5,4);
				\fill[DarkGray] (19,2) rectangle (20,4);

				\fill[RedOrange] (0,1.5) rectangle (1,2);
				\fill[RedOrange] (8.5,1.5) rectangle (11.5,2);
				\fill[RedOrange] (19,1.5) rectangle (20,2);

				\node at (0.5,1.75) {p+};
				\node at (9.5,1.75) {p+};
				\node at (19.5,1.75) {p+};

				%%% nmos:
				% body
				\fill[RedOrange] (17,1) rectangle (18.5,2);
				\node at (18,1.5) {p+};
				% source
				\fill[ProcessBlue] (15,1) rectangle (16.5,2);
				\node at (16,1.5) {n+};
				% drain
				\fill[ProcessBlue] (12,1) rectangle (13.5,2);
				\node at (13,1.5) {n+};

				%% gate:
				% gate oxide
				\fill[LightGray] (13.3,2) rectangle (15.2,2.1);
				% gate poly
				\fill[BrickRed] (13.3,2.1) rectangle (15.2,2.2);

				% metals
				\fill[DarkDarkGray] (17.5,2) rectangle (18,3);
				\fill[DarkDarkGray] (15.5,2) rectangle (16,3);
				\fill[DarkDarkGray] (15.5,3) rectangle (19,3.2); % connection pad GND

				\fill[DarkDarkGray] (14,2.2) rectangle (14.5,3);
				\fill[DarkDarkGray] (13.8,3) rectangle (14.7,3.2); % connection pad VG

				\fill[DarkDarkGray] (12.5,2) rectangle (13,3);
				\fill[DarkDarkGray] (12.3,3) rectangle (13.2,3.2); % connection pad VDD

				\fill[DarkDarkGray] (2,2) rectangle (2.5,3);
				\fill[DarkDarkGray] (4,2) rectangle (4.5,3);
				\fill[DarkDarkGray] (1,3) rectangle (4.5,3.2); % connection pad GND

				\fill[DarkDarkGray] (5.5,2.2) rectangle (6,3);
				\fill[DarkDarkGray] (5.3,3) rectangle (6.2,3.2); % connection pad VG

				\fill[DarkDarkGray] (7,2) rectangle (7.5,3);
				\fill[DarkDarkGray] (6.8,3) rectangle (7.7,3.2); % connection pad VDD

				% isolation oxides:
				\fill[NormalGray] (18,2) rectangle (19,3);
				\fill[NormalGray] (16,2) rectangle (17.5,3);
				\fill[NormalGray] (15.2,2) rectangle (15.5,3);
				\fill[NormalGray] (14.5,2.2) rectangle (15.2,3);
				\fill[NormalGray] (13.3,2.2) rectangle (14,3);
				\fill[NormalGray] (13,2) rectangle (13.3,3);
				\fill[NormalGray] (11.5,2) rectangle (12.5,3);

				\node at (1.5,3.5) {VDD};
				\node at (16.5,3.5) {Ground};
				\node at (5.5,3.4) {Input};
				\node at (14,3.4) {Input};
				\node at (7,3.4) {Output};
				\node at (12.5,3.4) {Output};

				\node at (0.5,4.2) {Field oxide};
				\node at (9.5,4.2) {Field oxide};
				\node at (19.5,4.2) {Field oxide};
			\end{tikzpicture}
		\end{center}
		\caption{Sectional view of a NMOS-PMOS transistor circuit}
	\end{figure}
\end{center}

\newpage
\section{Physics}
In this chapter we deal with all the physics related to solid state device manufacturing.
In case there is anything unclear, please look up this chapter and its sub-chapters.
\subsection{Getting doping from resistance}
In many cases the supplier will only provide the resistance per length specification for their substrate and won't give you the dopant concentration numbers.
In this case you will have to find these numbers out yourself by converting it from the numbers they've provided.

\begin{figure}[H]
	\centering
	\includegraphics[width=0.5\textwidth]{resistance_doping.png}
	\caption{R-L-dopant relation}
	\label{r_l_nd_relation}
\end{figure}

You can either use the graphics from \autoref{r_l_nd_relation} and determine the dopant concentration graphically, which is very very imprecise or use a online tool like the one from Solecon\footnote{\url{http://www.solecon.com/sra/rho2ccal.htm}}

Germanium is being included in this graphics just in case someone is going to fork this process based on Germanium substrate.
\newpage
\subsection{Infusion}
The redistribution process depends on the ratio of the solubility of the doping material in silicon and SiO$ _2$. At the Si/SiO$ _2$ interface the dopants are redistributed by segregation until the ratio of their concentration at the interface is the same as the ratio of their solubility in both materials. The ratio of dopant solubility is expressed by the segregation coefficient $m$ which is

\begin{equation}
\displaystyle m = \frac{
	\mathrm{solubility\ in\ silicon}
	}{
	\mathrm{solubility\ in\ SiO_2}
	}
\end{equation}

As listed in \autoref{table_diffusion_coeff} below there are dopant species which solubilize better in SiO$ _2$ than in silicon ($ m < 1$) and species which have a reversed behavior ($ m > 1$).
In case of $ m < 1$, as for Boron, the dopant concentration is enhanced at the SiO$ _2$ side, whereas beneath the interface, there is a dopant depletion at the silicon surface.
For reversed solubility ratios ($ m > 1$, like Phosphorus), only few dopant atoms penetrate the interface.
In order to obtain the by $ m$ determined concentration ratio at the interface, dopant atoms from deeper silicon zones diffuse back to the surface zone.
Therefore, the dopant concentration at the silicon surface is enhanced, as illustrated in \autoref{graph_figure}b.
In \autoref{graph_figure} $ C_c$ denotes the dopant concentration in the silicon surface zone before oxidation. $ x$ is the distance from the silicon surface.

\begin{center}
	\begin{figure}[h]
		\begin{center}
			\includegraphics[width=0.75\linewidth]{img349.png}
		\end{center}
		\caption{Schematic illustration of dopant redistribution}
		\label{graph_figure}
	\end{figure}
	\begin{table}[h]
		\begin{tabular}{|c|c|c|c|c|c|}
			\hline
			Dopant species &
			Boron &
			Phosphor &
			Antimon &
			Arsen &
			Gallium \\
			\hline
			$m$ &
			0.1-0.3 &
			10 &
			10 &
			10 &
			20 \\
			\hline
		\end{tabular}
		\caption{Segregation coefficients $m$ for important dopant species in silicon}
		\label{table_diffusion_coeff}
	\end{table}
\end{center}
\newpage
\section{Constant source diffusion (Predeposition)}

Although the diffusion process of donors and acceptors into the silicon crystal is a three dimensional process for simplicity we first only discuss the one dimensional mathematics for it in order to get a "simple" equation for the depth-time-temperature relation.

\begin{mdframed}[linewidth=2pt,linecolor=red]
This is only valid for a constant source of dopants on the surface of the wafer (gas, for instance).
These equations are used for predicting the pre-deposition step (in case this process would be adapted by someone for predeposition instead of ion implant)
\end{mdframed}

We start with Ficks\footnote{\url{https://en.wikipedia.org/wiki/Fick\%27s_laws_of_diffusion}} law (for all German speakers: Yes that's his name) where the dopant concentration N is coupled with time and place
\begin{equation}
\frac{\partial N}{\partial t} = D \cdot \frac{\partial^2 N}{\partial x^2}
\end{equation}

The diffusion coefficient is as well material as well as temperature dependent  and can be calculated with the following equation:
\begin{equation}
D = D_0 \cdot \exp\left(-\frac{E_a}{k \cdot T}\right)
\end{equation}
With $k=8.62 \cdot 10^{-5} \frac{eV}{K}$ being the Boltzman constant and in table \ref{table:absolute_diffusion_coefficients} we can see the $D_0$ and $E_a$ values for the most common materials\footnote{ISBN 3-8023-1588:Hoppe Bernhard, Mikroelektronik 2, Page 24, Table 2.1} which we can use within the further calculations for our well dimensioning phases. The temperature usually is in the area of $1000\degree C$ or $1273.15 K$.
\begin{table}[H]
	\label{table:absolute_diffusion_coefficients}
	\centering
	\begin{tabular}{|c|c|c|}
		\hline
		Element &
		$D_0$ $\left[\frac{cm^2}{s}\right]$ &
		$E_a$ $\left[eV\right]$ \\
		\hline
		P &
		10.50 &
		3.69 \\
		\hline
		As &
		0.32 &
		3.56 \\
		\hline
		Sb &
		5.60 &
		3.95 \\
		\hline
		B &
		10.50 &
		3.69 \\
		\hline
		Al &
		8.00 &
		3.47 \\
		\hline
		Ga &
		3.60 &
		3.51 \\
		\hline
		Cu &
		0.0025 &
		0.65 \\
		\hline
	\end{tabular}
	\caption{$D_0$ and $E_a$ values for boron and phosphorus}
\end{table}

The law stated above
\begin{equation}
\frac{\partial N}{\partial t} = D \cdot \frac{\partial^2 N}{\partial x^2} 
\end{equation}
has the same form as the temperature conductivity equation (Laplace) for which we already have a general solution
\begin{equation}
\frac{\partial u}{\partial t} = a^2 \cdot \frac{\partial^2 u}{\partial x^2} 
\end{equation}

Which means that we can map the general solution for the temperature conductivity equations after Laplace
\begin{equation}
u(x,t) = \frac{1}{2 \cdot a \cdot \sqrt{\pi \cdot t}} \cdot \int_{-\infty}^{\infty}{f(a)\cdot\exp\left(\frac{-(x-a)^2}{4 \cdot a^2 \cdot t^2}\right)}da
\end{equation}
to our Ficks law with $a=\sqrt{D}$ and $u=N$
\begin{equation}
N(x,t) = \frac{1}{2 \cdot \sqrt{D} \cdot \sqrt{\pi \cdot t}} \cdot \int_{-\infty}^{\infty}{f(\sqrt{D})\cdot\exp\left(\frac{-(x-\sqrt{D})^2}{4 \cdot D \cdot t^2}\right)}da
\label{eq:ficks_law_expanded}
\end{equation}
with the edge conditions
\begin{equation}
N( x=0 , t > 0 ) = N_0
\end{equation}
\begin{equation}
N( x \geq 0 ,  t = 0 ) = 0
\end{equation}
we get the resulting function from the solving process for the Laplace temperature conduction equations
\begin{equation}
u(x,t)=u_0 \cdot\erfc\left(\frac{x}{2 \cdot a \cdot \sqrt{t}}\right)
\end{equation}
with the error function being an integral of the form
\begin{equation}
\erfc(z)
=
\left(1-\frac{2}{\sqrt{\pi}}\right)\cdot\int_0^z{e^{-a^2}}da
\end{equation}

Or in case of our dopant concentration equation we can replace a with the square root of the diffusion coefficient in order to get the error function for our dopant density equation:
\begin{equation}
\erfc(z)
=
\left(1-\frac{2}{\sqrt{\pi}}\right)\cdot\int_0^z{e^{-D}}d\sqrt{D}
\end{equation}
\begin{equation}
N(x,t)
=
N_0 \cdot\erfc\left(\frac{x}{2 \cdot \sqrt{D \cdot t}}\right)
=
N_0 \cdot\erfc\left(\frac{x}{x_l(t)}\right)
\end{equation}
Now we can extract the layer thickness and the depth of the well in dependency of the time and the temperature, respectively:
\begin{equation}
x_l(t) = 2 \cdot \sqrt{D \cdot t}
\end{equation}
\begin{equation}
x_l(t) = 2 \cdot \sqrt{D_0 \cdot \exp\left(-\frac{E_a}{k \cdot T}\right) \cdot t}
\end{equation}

And plot the result for multiple different drive in times
\begin{figure}[H]
	\centering
	\includegraphics[scale=0.5]{dopants_depth.png}
	\caption{Different predeposition times}
\end{figure}

We can now describe the dosage based on the time and temperature of the diffusion

\begin{equation}
Q=\frac{2}{\sqrt{\pi}} \cdot N_0 \cdot \sqrt{D \cdot t}
\end{equation}

Where $N_0$ (concentration at the surface) equals the maximum solubility of a given element (e.g. boron) within the given medium (e.g. silicon).\footnote{If someone really wants to do this in his basement he can google these values and make a pull request}

\newpage
\subsection{Ion implant and drive-in}
We can use the following equation to calculate the carrier distribution after implantation and annealing
\begin{equation}
N(x,t)
=
\frac{Q}{\sqrt{2\pi} \sqrt{\Delta R_p^2+2 D t}}
\cdot
\exp\left(
-\frac{(x-R_p)^2}{2 \cdot (\Delta R_p^2+2 D t)}
\right)
\end{equation}

 $R_p$
\footnote{\url{http://cleanroom.byu.edu/rangestraggle}}
\subsection{Drive-in (limited source diffusion)}\label{physics_drive_in}
After pre-deposition or ion implant of the initial dosage we need to drive in the ions deeper into the crystal.
In order to prevent back-diffusion into the gas we seal off the oxide window with another layer of oxide in order to make sure that all the dopants stay inside the silicon crystal.

\begin{figure}[H]
	\centering
	\includegraphics[scale=0.5]{dopants_drive_in_depth.png}
	\caption{Drive-in well depths and concentrations}
\end{figure}

We set the condition that the pre-deposition/implant depth is much lower than the depth of the final diffused volume with the following inequation:
\begin{equation}
x_e = 2 \cdot \sqrt{D_e \cdot t_e} \gg 2 \cdot \sqrt{D_v \cdot t_v} = x_v
\end{equation}
Where $x_v$ is the the depth of the predeposition/implant step.

By neglecting the distribution thickness of the original implantation dosage and assuming that it's comparably thin compared to the medium thickness we can replace  $f(a) \approx \delta(a)$ within \autoref{eq:ficks_law_expanded} which makes
\begin{equation}
N(x,t) = \frac{1}{2 \cdot \sqrt{D} \cdot \sqrt{\pi \cdot t}} \cdot \int_{-\infty}^{\infty}{f(\sqrt{D})\cdot\exp\left(\frac{-(x-\sqrt{D})^2}{4 \cdot D \cdot t}\right)}da
\end{equation}
become
\begin{equation}
N(x,t)
=
\frac{1}{2 \cdot \sqrt{D} \cdot \sqrt{\pi \cdot t}} \cdot \int_{-\infty}^{\infty}{\delta(\sqrt{D})\cdot\exp\left(\frac{-(x-\sqrt{D})^2}{4 \cdot D \cdot t}\right)}da
\end{equation}
and finally
\begin{equation}
N(x,t)
=
\frac{Q}{\sqrt{\pi \cdot D_e \cdot t}} \cdot \exp\left(\frac{-x^2}{4 \cdot D_e \cdot t}\right)
\end{equation}
\subsection{Vertical diffusion and junction formation (Well formation)}\label{building_wells}
The goal of most diffusions is to form pn junctions by converting p-type material to n-type material or vice versa.
In \autoref{junction_formation_plot}, for example, the wafer is uniformly doped n-type material with a concentration indicated by $N_B$, and the diffusing impurity is boron.
The point at which the diffused impurity profile intersects the background concentration is the mettalurgical junction depth ($x_j$).
The net impurity concentration at $x_j$ is zero.
Setting $N(x)$ equal to the background concentration $N_B$ at $x=x_j$ yields\footnote{Gerold W. Neudeck and Robert F. Pierret, Modular series on solid state devices, Volume V, Chapter 4}
\begin{equation}
x_j
=
2 \cdot\sqrt{D \cdot t \cdot \ln\left(\frac{N_0}{N_B}\right)}
\end{equation}
and
\begin{equation}
x_j
=
2 \cdot\sqrt{D \cdot t}
\cdot
\erfc^{-1}\left(\frac{N_B}{N_0}\right)
\end{equation}

for the Gaussian and complementary error function distributions, respectively.

In \autoref{junction_formation_plot}, the boron concentration $N$ exceeds $N_B$ to the left of the junction, and this region is p-type.
To the right of $x_j$, $N$ is less than $N_B$, and this region remains n-type.

To calculate the junction depth, we must know the background concentration $N_B$ of the original wafer.
Look at \autoref{r_l_nd_relation} for this purpose.

\begin{figure}[H]
	\centering
	\includegraphics[scale=0.5]{well_formation1.png}
	\includegraphics[scale=0.5]{well_formation2.png}
	\caption{Formation of a pn junction by diffusion: (a) An example of a p-type Gaussian diffusion into a uniformly doped n-type wafer; (b) net impurity concentration in the wafer.}
	\label{junction_formation_plot}
\end{figure}


\newpage
\section{MOS Capacitance}
\url{https://ecee.colorado.edu/~bart/book/book/chapter6/ch6_3.htm}

\subsection{Threshold voltage ($V_T$)}
The formula for calculating the threshold voltage of a MOS device is the following:
\begin{equation}
V_T = V_{t-mos} + V_{FB}
\end{equation}
where $V_{t-mos}$ is the ideal threshold voltage of an ideal MOS capacitor and $V_{FB}$ is what is termed flat-band voltage and $V_{t-mos}$ is the threshold.
The MOS threshold voltage, $V_{t-mos}$ is calculated by considering the MOS capacitor structure that form the gate of the MOS transistor.
The ideal threshold voltage may be expressed as:
\begin{equation}
V_{t-mos}=2 \phi_F + \frac{Q_b}{C_{ox}}
\end{equation}
with
\begin{equation}
Q_b
=
\sqrt{2 \epsilon_{Si} q N_A ( 2 \phi_F + V_{SB}) }
\end{equation}
\begin{equation}
\phi_F
=
V_{th} ln\left(\frac{N_A}{N_i}\right)
=
\frac{k T}{q} ln\left(\frac{N_A}{N_i}\right)
\end{equation}
where $C_{ox}$ is the oxide capacitance and $Q_b$ which is called the bulk charge term and $V_{th}$ is the thermal voltage.\footnote{\url{https://en.wikipedia.org/wiki/Boltzmann_constant\#Role_in_semiconductor_physics:_the_thermal_voltage}}

Because we put source and bulk onto the same potential in our CMOS circuit as shown in \autoref{process_design} we can put $V_{BS}=0$ which results in the following new equation for $Q_B$:
\begin{equation}
V_{BS}=0
\Rightarrow
Q_b
=
\sqrt{2 \epsilon_{Si} q N_A 2 \phi_F }
=
2 \sqrt{\epsilon_{Si} q N_A \phi_F }
\end{equation}

 $V_{FB}$, is given by:
\begin{equation}
V_{FB}
=
\phi_{MS}-\frac{Q_f}{C_{ox}}-\frac{1}{C_{ox}}\int_{0}^{t_{ox}}\frac{x}{x_{ox}}\rho(x) dx
\end{equation}
We assume the oxide doesn't contain any charge, which makes the equation fall together into
\begin{equation}
\rho(x)=0
\Rightarrow
V_{FB}
=
\phi_{MS}-\frac{Q_f}{C_{ox}}
\end{equation}
with
\begin{equation}
\phi_{MS}
=
\phi_{M} - \phi_{S}
=
\phi_{M} -  \left( \frac{E_g}{2 q} + \phi_F \right)
=
\phi_{M} - \frac{E_g}{2 q} - \phi_F
\end{equation}

We now can put the equation together:
\begin{equation}
V_T = V_{t-mos} + V_{FB}
\end{equation}
\begin{equation}
V_{t-mos}=2 \phi_F + \frac{Q_b}{C_{ox}}
\Rightarrow
V_T = 2 \phi_F + \frac{Q_b}{C_{ox}} + V_{FB}
\end{equation}
\begin{equation}
V_{FB}
=
\phi_{MS}-\frac{Q_f}{C_{ox}}
\Rightarrow
V_T = 2 \phi_F + \frac{Q_b}{C_{ox}} + \phi_{MS}-\frac{Q_f}{C_{ox}}
\end{equation}
\begin{equation}
Q_b
=
2 \sqrt{\epsilon_{Si} q N_A \phi_F }
\Rightarrow
V_T = 2 \phi_F + \frac{2 \sqrt{\epsilon_{Si} q N_A \phi_F }}{C_{ox}} + \phi_{MS}-\frac{Q_f}{C_{ox}}
\end{equation}


\begin{equation}
\phi_{MS}
=
\phi_{M} - \frac{E_g}{2 q}-\phi_F
\Rightarrow
V_T = 2 \phi_F + \frac{2 \sqrt{\epsilon_{Si} q N_A \phi_F }}{C_{ox}} + \phi_{M} - \frac{E_g}{2 q}-\phi_F -\frac{Q_f}{C_{ox}}
\end{equation}
\begin{equation}
\Rightarrow
V_T = \phi_F + \frac{2 \sqrt{\epsilon_{Si} q N_A \phi_F }-Q_f}{C_{ox}} + \phi_{M} - \frac{E_g}{2 q}
\end{equation}


\begin{equation}
\phi_F
=
\frac{k T}{q} ln\left(\frac{N_A}{N_i}\right)
\end{equation}
\begin{equation}
\Rightarrow
V_T = \frac{k T}{q} ln\left(\frac{N_A}{N_i}\right) + \frac{2 \sqrt{\epsilon_{Si} q N_A \frac{k T}{q} ln\left(\frac{N_A}{N_i}\right) }-Q_f}{C_{ox}} + \phi_{M} - \frac{E_g}{2 q}
\end{equation}
\begin{equation}
 \boxed{
\Rightarrow
V_T = \frac{k T}{q} ln\left(\frac{N_A}{N_i}\right) + \frac{2 \sqrt{\epsilon_{Si} N_A k T ln\left(\frac{N_A}{N_i}\right) }-Q_f}{C_{ox}} + \phi_{M} - \frac{E_g}{2 q}
}
\end{equation}


With the variables being the following:
\begin{itemize}
\item $N_A$ is the substrate doping (which in the case of our chosen basis substrate is approximately $9\times10^{14}cm^{-3}$ (see \autoref{process_overview}))
\item $N_i$ is the carrier concentration in intrinsic (undoped) silicon. $N_i$ is equal to $1.45 \times 10^{10} cm^{-3}$ at 300\degree K
\item $\phi_M = 4.1 V$ is the "work function" of our metal at the gate (multiplied with q) (Only Aluminum is practical for this process)
\item $E_g$ is the band gap energy of silicon: $\left(1.16-0.704 \times 10^{-3} \frac{T^2}{T+1108} \right)^5$
\item $C_{ox}$ is the oxide capacitance
\end{itemize}


\section{Threshold voltage ($V_T$) adjustment}
At some point in the future this will be of very high relevance, because the lower the size of the transistors becomes, the higher the offset to $V_{Tp}$ and $V_{Tn}$ needs to be in order to stay on TTL 5V logic level, or at least compensate for the lowered voltages in order to reach the 3.3V CMOS logic levels.

Adjustment of the threshold voltage can be achieved by:
\begin{itemize}
\item A relatively small dose $N_I$ (units: ions/$cm^2$) of dopant atoms is implanted into the near-surface  region of the semiconductor.
\item When the MOS device is biased in depletion or inversion, the implanted dopants add to (or substract from) the depletion charge near the oxide-semiconductor interface
\end{itemize}

The formula to calculate the voltage offset is:
\begin{equation}
\Delta V_T = -\frac{q N_I}{C_{ox}} 
\left\{\begin{matrix}
N_I > 0\ for\ donor\ atoms\ (Phosporus/N) \\
N_I < 0\ for\ acceptor\ atoms\ (Boron/P)
\end{matrix}\right.
\end{equation}

\newpage
\section{Chemistry}
\subsection{Etching}
\subsubsection{Silicon dioxide}
A very "selective" chemical for SiO2 - i.e. does not etch silicon at all - is hydrofluoric acid (HF). If used directly such etchant has a too fast and aggresive action on the oxide, making very difficult the undercut and the linewidth control. For such reason, HF is universally used as a "buffered" solution, which can keep the etch rate low and constant, by moderating the PH level of the bath. This allows the etching time to be reliably correlated to the etching depth.

The industry standard buffered hydrofluoric acid solution (BHF) has the following formulation:

- 6 volumes of ammonium floride (NH4F, 40% solution) - 1 volume of HF.

This can be prepared, for example, by mixing 113 g of NH4F in 170 ml of H2O, and adding 28 ml of HF. The etch rate at room temperature can range from 1000 to 2500 Å/min. This depends on the actual density of the oxide which, as an amorphous layer, can have a more compact structure (if thermally grown in is oxygen) or less compact (if grown by CVD). The following etching reaction holds:

\begin{equation}
SiO_2 + 6HF \rightarrow H_2SiF_6 + H_2O
\end{equation}

where H2SiF6 is water soluble.

Sometimes the BHF is prepared and stored according to the formulation given above, but it is 7:1 diluited in water just before being used. This allows an even better control of the etching rate. It is good practice to use the diluited BHF only once and then discard it, in order to assure process repeatability. If a diluited BHF bath at 35°C is used, the etching rate for thermal oxide is around 800Å/min.

Another popular etching formulation is the P-etch:

60 volumes of H2O + 3 vol. of HF + 2 vol. of HNO3, that is: 300 ml of H2O + 15 ml of HF + 10 ml of HNO3.

The P-etch action is strongly dependent on oxide density, as it results from the growth technique. An example is reported in the literature\footnote{A. Pliskin, J.Vac.Sci Technol., vol. 14, p.1064, 1977}, indicating 120 Å/min for thermal oxide and 250-700 Å/min for sputtered oxide.

A slow etching bath is preferred for opening mask windows for a silicon substrate. However, the etching process could be used just for removing the oxide film from the whole surface. In this case the etching speed is not critical, and a fast solution can be used, such as HF diluited 1:10 in water. The etching time can be easily evaluated by visually inspecting the surface. Once the oxide film is removed, the metal-grey color of the silicon surface appears.

Sometimes a very light etch is required, for removing just a few atomic layers. This is the case of surface cleaning and decontamination. HF diluited 1 : 50 in water can be used. The etching speed will be around 70 Å / min. For example, a typical 50 Å "native" oxide on silicon can be removed with a 45 - 50 sec light etch.


\subsection{Growing silicon nitride}\label{chemistry_growing_nitride}
In order to grow a high quality layer of silicon nitride on top of a silicon wafer which is adapted to be patterned and to serve as a mask for diffusion or implantation of selected impurities, the wafer is being best put into a chamber evacuated to a pressure less than about 1 Torr and heated to between 650 and 900 \degree C;
A gaseous mixture comprising primarily ammonia and a silicon compound is being flooded into that chamber with a silicon compound flow rate of greater than approximately 12 cubic centimeters per minute and the ammonia and selected silicon compound having a ratio of relative concentrations in the range of 4:1 and 20:1.\footnote{\url{http://www.freepatentsonline.com/4395438.html}}
The growth rate will be around 50 Angstroms per minute.
That setup is called Low-Pressure Chemical Vapor Deposition (LPCVD), which is commonly available in basically any semiconductor manufacturing plant or laboratory.

\newpage
\section{Process design}
We need to optimize our process to be TTL compatible (5V logic levels) and at the same time being as fast and power efficient as possible.
In order to have a good propagation delay with a technology node of around $1\mu m$ we will have to have gates with up to four stacked MOS transistors.

Acceptable input signal voltages range from 0 volts to 0.8 volts for a low logic state, and 2 volts to 5 volts for a high logic state.
Acceptable output signal voltages shall range from 0 volts to 0.5 volts for a low logic state, and 2.7 volts to 5 volts for a high logic state\footnote{\url{https://www.allaboutcircuits.com/textbook/digital/chpt-3/logic-signal-voltage-levels}}

\begin{figure}[H]
	\centering
	\includegraphics[scale=0.5]{logic_levels.png}
	\caption{TTL logic levels}
	\label{TTL_logic_levels}
\end{figure}

As shown in \autoref{TTL_logic_levels} we have some margin to make our PMOS and NMOS transistors work with each other in order to form a CMOS circuit which is actually working without getting warm.

Or more clearly defined
\begin{equation}
V_{off} \leq 0.8V
\end{equation}
and
\begin{equation}
V_{on} \geq 2V
\end{equation}
which are limits, elementary to our design.

\begin{figure}[H]
	\centering
	\includegraphics[scale=0.5]{cmos_3v3.png}
	\caption{CMOS 3.3V logic levels}
	\label{CMOS_logic_levels}
\end{figure}

This means that we also will be compatible to CMOS logic level output pins since their ON/OFF levels are within our tolerance range\footnote{\url{https://learn.sparkfun.com/tutorials/logic-levels/33-v-cmos-logic-levels}} as it is shown in \autoref{CMOS_logic_levels}.
\newpage
\input{process_design_requirements.tex}
\section{Substrate}

The Hong University of science and technology (short HKUST) provides us with two types of wafers.

\begin{itemize}
	\item Prime Grade Silicon Wafer, [100] N-type
	\begin{itemize}
		\item Front-side polished, backside etched
		\item Dopant: Phosphorus
		\item Thickness: $525 \mu m \pm 25 \mu m$
		\item Resistivity:4 to 7ohm-cm
		\item Growth Method: CZ
		\item Diameter: 100mm +/- 0.5 mm
		\item Primary  \& secondary flat locations: (In compliance with the SEMI)
			\begin{itemize}
				\item Carbon concentration<2.5 x e16 atm/cc
				\item Oxygen concentration < $9.0 \cdot 10^{17} \frac{atm}{cc}$
				\item $TTV < 10 \mu m$
				\item $TIR < 6 \mu m$
				\item $Bow/Warp < 40 \mu m$
			\end{itemize}
	\end{itemize}
	\item Prime Grade Silicon Wafer, [100] P-type
	\begin{itemize}
		\item Front-side polished, backside etched
		\item Dopant: Boron
		\item Thickness: $525 \mu m \pm 25 \mu m$
		\item Resistivity:15 to 25 ohm-cm
		\item Growth Method: CZ
		\item Diameter: 100mm +/- 0.5 mm
		\item Primary  \& secondary flat locations: (In compliance with the SEMI)
			\begin{itemize}
				\item Carbon concentration<2.5 x e16 atm/cc
				\item Oxygen concentration < $9.0 \cdot 10^{17} \frac{atm}{cc}$
				\item $TTV < 10 \mu m$
				\item $TIR < 6 \mu m$
				\item $Bow/Warp < 40 \mu m$
			\end{itemize}
	\end{itemize}
\end{itemize}

For this process the p-doped mono crystalline silicon substrate is being used, but forks and modifications will be very well possible based on a Graphene substrate or alike, still under the LSPL.
The starting material is a p-doped <100> oriented mono crystalline silicon wafer\\

\textbf{Reasons for using p-doped substrate}:\begin{itemize}
\item We can't use two different substrates for our design because in the design both PMOS and NMOS is present.
We have to choose which is more beneficial from fabrication point of view.
In general or say it's true that NMOS devices are always more in the Semiconductor Industry in comparison to PMOS devices.
For your reference-SRAM requires 6 transistors (4 NMOS, 2 PMOS).
\item Another reason for more number of NMOS is because of difference of mobility of electron and holes.
Electron mobility is almost twice of holes mobility and because of this ON-RESISTANCE of n-channel device is half of p-channel device with the same geometry and under the same operating conditions.
That means to achieve same impedance size of n-channel transistors is almost half of p-channel devices.
Same thing I can say in the different way that for same size of wafer, we can have more number of NMOS (means can perform more logical operation) in comparison to PMOS.
\item Since we only have the choice between P and N doped substrate, we use P doped substrate, because of the carrier mobility
\end{itemize}

Using the method from \autoref{r_l_nd_relation} we get a doping concentration between $8.76 \cdot 10^{14} \frac{1}{cm^3}$ and $5.23 \cdot 10^{14} \frac{1}{cm^3}$.
The average of this range is $N_B = \frac{8.76+5.23 }{2} \cdot 10^{14} \frac{1}{cm^3} \approx 7 \cdot 10^{14} \frac{1}{cm^3}$
\subsection{Isolation}
For the isolation (\autoref{sti})  in this design the STI approach is being chosen.
Shallow trench isolation (STI), also known as box isolation technique, is an integrated circuit feature which prevents electric current leakage between adjacent semiconductor device components.\footnote{\url{https://www.google.com/patents/US7985656}}
STI is generally used on CMOS process technology nodes of 250 nanometers and smaller.\\

\textbf{Reasons for using box isolation}:\begin{itemize}
\item We want to be forward compatible to future LibreSilicon nodes with a size of 100nm or smaller
\item It simplifies the construction of the gate and allows us to use Aluminum instead of Polysilicon for the gate contact
\end{itemize}

\subsection{Interconnect}
The interconnects and the gate electrode are being made using Aluminum which is a very commonly used material to do interconnects in low-frequency and low-resolution applications \\

\textbf{Reasons for using Aluminum}:\begin{itemize}
\item It's a well explored material for interconnect with a lot of literature on how to process it
\item Aluminum is easy to etch compared to copper
\item It isn't contaminating everything like copper does and doesn't require special separated setup for handling
\item The machines at HKUST can do CMP for copper only on 4 inch wafers which would limit us in wafer size
\end{itemize}

\begin{mdframed}[linewidth=2pt,linecolor=green]
As soon as we've got CMOS all figured out, we will tackle copper interconnect in release 2.0
\end{mdframed}
\newpage
\subsection{MOS gate material}


\subsection{MOS gate thickness}\label{gate_dimensioning}
As the continuous down-scaling of the device size has lead to very thin gate oxides, the leakage current that can flow from the channel to the gate comes into the order of the subthreshold leakage current and the gate cannot be considered as an ideally insulated electrode anymore.
This affects the circuit functionality and increases the standby power consumption due to the static gate current.
For dynamic logic concepts the gate leakage drastically reduces the maximum clock cycle time\footnote{N. Wang, Digital MOS Integrated Circuits, Prentice-Hall, Englewood Cliffs, NJ, 1989}.
Two tunneling mechanisms are responsible for the gate leakage, Fowler-Nordheim tunneling and direct tunneling\footnote{A. Schenk and G. Heiser, "Modeling and Simulation of Tunneling through Ultra-Thin Gate Dielectrics"J.Appl.Phys., vol. 81, no. 12, pp. 7900, 1997}.
The gate leakage increases exponentially as the oxide thickness is reduced.
This limits the down-scaling of the oxide thickness to about 1.5-2 nm when looking at the total standby power consumption of a chip\footnote{Y. Taur, "The Incredible Shrinking Transistor," IEEE Spectrum, pp. 25-29, July 1999.}.
To further decrease the effective oxide thickness alternative high dielectric constant materials can be used\footnote{S. Thompson, P. Packan, and M. Bohr, "MOS Scaling: Transistor Challenges for the 21st Century," Intel Technology Journal, vol. Q3, 1998}.
On the other hand, a thin gate oxide reduces the short-channel effect and improves the driving capabilities of a MOS transistor.
However, a tradeoff between this benefit and the gate leakage is necessary.\\

With $1 \mu m$ we don't have to worry about this leakage yet because our gate oxide thickness is too high for these effects to actually become a problem, but we want to do our home work already in preparation of scale-down and also for curiosity.

We for now just use 40 nm. That's still doable with a precision high enough when using dry oxidation and a temperature of 1000\degree Celsius.

\subsubsection{Subthreshold leakage}\label{sub_threshold_leakage}

The sub-threshold leakage current can be calculated with\footnote{\url{http://ecee.colorado.edu/\~bart/book/book/chapter3/ch3_4.htm\#3_4_2}}
\begin{equation}
I_{sub}
=
I_0
\cdot
\left(1-\exp\left(-\frac{V_{ds}}{V_{th}}\right)\right)
\cdot
\exp\left(\frac{V_{gs}-V_{T}}{n \cdot V_{th}}\right)
\end{equation}

where
\begin{equation}
I_0 = \frac{W}{L} \mu_0 V_{th}^2 \sqrt{\frac{N_A \cdot q \cdot \epsilon_{Si}}{2 \cdot \phi_{sub}}}
\end{equation}
$V_{th}=26mV$ is the thermal voltage, $V_T$ is the threshold voltage, $V_{ds}$ and $V_{gs}$ are the drain-to-source and gate-to-source voltages respectively.
W and L are the effective transistor width and length, respectively. $C_{ox}$ is the gate oxide capacitance, $\mu_0$ is the carrier mobility and $n=1+\frac{C_{dep}}{C_{ox}}$ is the subthreshold swing coefficient.

First of all, lets say $W=L$ which leads to a square:
\begin{equation}
I_0 = \mu_0 V_{th}^2 \sqrt{\frac{N_A \cdot q \cdot \epsilon_{Si}}{2 \cdot \phi_{sub}}}
\end{equation}

With
\begin{itemize}
\item $\epsilon_0 = 8.85 \cdot 10^{-14}\frac{F}{cm}. $ is the electric permittivity in vacuum
\item $\epsilon_{ox} =3.9 \cdot \epsilon_0$ is the relative permittivity of silicon dioxide
\item $\epsilon_{Si} =11.68 \cdot \epsilon_0$ is the relative permittivity of silicon
\end{itemize}

The carrier mobility $ \mu_0$ can be calculated with\footnote{\url{https://ecee.colorado.edu/\~bart/book/book/chapter2/ch2_7.htm\#2_7_2}}
\begin{equation}
 \mu(N) =  \mu_{min} + \frac{ \mu_{max}- \mu_{min}}{1+\left(\frac{N}{N_r}\right)^\alpha}
\end{equation}
using the fitting parameters from \autoref{fitting_parameters_mobility}

\begin{table}[H]
	\centering
	\begin{tabular}{|c|c|c|c|}
		\hline
		{} &
		\textbf{Arsenic} &	
		\textbf{Phosphorus} &
		\textbf{Boron} \\
		\hline
		$\mu_{min} [\frac{cm^2}{Vs}]$ &
		52.2 &
		68.5 &		
		44.9 \\
		\hline
		$\mu_{max} [\frac{cm^2}{Vs}]$ &
		1417 &
		1414 &
		470.5 \\
		\hline
		$N_r [\frac{1}{cm^3}]$ &
		$9.68 \cdot 10^{16}$ &
		$9.20 \cdot 10^{16}$ &
		$2.23 \cdot 10^{17}$ \\
		\hline
		$\alpha$ &
		0.68 &
		0.711 &
		0.719 \\
		\hline
	\end{tabular}
	\caption{Parameters for calculation of the mobility as a function of the doping density}
	\label{fitting_parameters_mobility}
\end{table}

\subsubsection{Gate tunneling current}

The tunneling of electrons (or holes) from the bulk and source/drain overlap region through the gate oxide potential barrier into the gate (or vice-versa) is referred as gate oxide tunneling current.
This phenomenon is related with the MOS capacitance concept.
There are three major gate leakage mechanisms in a MOS structure.
The first one is the electron conduction-band tunneling (ECB), where electrons tunneling from conduction band of the substrate to the conduction band of the gate (or vice versa).
The second one is the electron valence-band tunneling (EVB). In this case, electrons tunneling from the valence band of the substrate to the conduct band of the gate.
The last one is known as hole valence-band (HVB) tunneling, where holes tunneling from the valence band of the substrate to the valence band of the gate (or vice- versa)

Each mechanism is dominant or important in different regions of operation for NMOS and PMOS transistors. For each mechanism, gate leakage current can be modeled by
\begin{equation}
I = W \cdot L \cdot A \cdot \left(\frac{V_{ox}}{t_{ox}}\right)^2\exp\left(\frac{-B\left(1-\left(1-\frac{V_{ox}}{\phi_{ox}}\right)^{\frac{3}{2}}\right)}{\frac{T_{ox}}{t_{ox}}}\right)
\end{equation}

\newpage
\subsection{NMOS threshold}\label{nmos_gate_dimensioning}
First we take a look at the worst case of 4 stacked NMOS transistors, which is the highest stacking amount which will be possible in technologies relying on this process.

\input{process_design_graphics/schematic_AND4.tex}
As shown in  \autoref{TTL_logic_levels} our acceptable voltages for our CMOS "ON" state range from 2V to VDD (which typically is around $5V\pm0.25V$)

\begin{equation}
V_{on} \geq 2V
\end{equation}

Because there are four transistors dividing the voltage by being in series, the power supply voltage is being divided by the amount of transistors in series.
In order to match the threshold voltages of all of the transistors, which is needed for a working digital logic, the following equation need to be satisfied

\begin{equation}
V_{on} > 4 \cdot V_{Tn}
\end{equation}

Lets assume the worst case with

\begin{equation}
V_{on} = 2V
\end{equation}

Which leads to the required worst case threshold tolerance value

\begin{equation}
4 \cdot V_{Tn} < 2V
\Rightarrow
V_{Tn} < 500mV
\end{equation}
 
 Using the equation from \autoref{transistor_threshold_calculation} we can now set in this value and play around with the parameters.
 
 \begin{equation}
V_T= \frac{1.66 \cdot 10^{-7}}{C_{ox}}-0.094
\end{equation}

For safety reasons we take 10\% of the worst case as a error margin which leads us to the following new equation to solve
 \begin{equation}
V_T = 450mV
\end{equation}

This gives us the constraint of
\begin{equation}
0.45 V = \left( \frac{1.66 \cdot 10^{-7}}{C_{ox}}-0.094 \right) V
\end{equation}

We solve for the target capacity (in $\frac{F}{cm^2}$) which gives us the desired threshold voltage:
\begin{equation}
\Rightarrow
C_{ox}
=
3.057 \cdot {10}^{-7} \frac{F}{cm^2}
\end{equation}
\begin{equation}
\Rightarrow
C_{ox}
=
305.7 \frac{nF}{cm^2}
\end{equation}
\newpage
\subsection{PMOS threshold}\label{pmos_dimensioning}
Now we take a look at the worst case of 4 stacked PMOS transistors, which is the highest stacking amount which will be possible in technologies relying on this process.
\input{process_design_graphics/schematic_OR4.tex}


With our $\phi_F$ for a doped n-well
\begin{equation}
\phi_F
=
V_{th} ln\left(\frac{{N_A \cdot N_D}}{N_i^2}\right)
\end{equation}
\newpage
\subsection{MOS gate verification}\label{gate_verify}
Now we have to verify that we don have relevant leakage on the gate with the thresholds and doping values we've set.

For that we use the formulas from \autoref{sub_threshold_leakage}

We can now plot multiple leakages for N- and P-channel transistors with a gate oxide thickness\footnote{See simulation/gate.wxmx} and with a surface concentration of $1e16\frac{1}{cm^3}=1e22\frac{1}{m^3}$ each
\begin{figure}[H]
	\centering
	\includegraphics[width=\textwidth]{subthreshold_leagage.png}
	\caption{Subthreshold leakage plot(in Ampere)}
	\label{leagage_currents}
\end{figure}
In \autoref{leagage_currents} we see that with our gate oxide thickness this is really no problem, as we had expected.
From 0V up to 5V and further there is basically no leakage on the gate from the sub threshold current with $V_{Tn} \approx 0.39V$ and $V_{Tp} \approx -0.30V$.
That's good enough, as we will see in \autoref{nmos_dimensioning} and \autoref{pmos_dimensioning}.
There is actually a reduction of current when reaching the threshold because of the inversion of the capacity in the depletion zone\footnote{\url{https://people.eecs.berkeley.edu/~hu/Chenming-Hu_ch5.pdf}}, but I didn't include this into the calculation, because "TL;DR".
It's a TODO for release 2.1 of this process which will go sub $1 \mu m$
\newpage

\section{Process steps}
\section{Process}

Below the general flow chart of the overall process flow can be seen in \autoref{full_flow}.
These process steps will be discussed within the following sections.

\begin{center}
	\begin{figure}[h]
		\begin{center}
			\begin{tikzpicture}[node distance = 3cm, auto, thick,scale=0.5, every node/.style={transform shape}]
				% Place nodes
				\node [block] (clean1) at (0,15) {Initial cleaning};
				\node [block, right of=clean1] (well) {Well};
				\node [block, below of=well] (fox) {Channel stop and field oxide};
				\node [block, below of=fox] (active) {Active};
				\node [block, below of=active] (pp) {p+ Implant};
				\node [block, below of=pp] (np) {n+ Implant};
				\node [block, right of=np] (metal) {Metal 1-6};
				% Draw edges
				\path [line] (clean1) -- (well);
				\path [line] (well) -- (fox);
				\path [line] (fox) -- (active);
				\path [line] (active) -- (pp);
				\path [line] (pp) -- (np);
				\path [line] (np) -- (metal);
				
				\draw[dotted] (1.5,2) rectangle (4.5,16);
				\node at (6,8) {CMOS process};
			\end{tikzpicture}
		\end{center}
		\caption{Frontend and backend process flow}
		\label{full_flow}
	\end{figure}
\end{center}

The five starting overall process steps are part of an overall active part of the technology, while the final metal (respectively contact) layers will be used for making a contact between the logic gates and macro cells and making them available to the exterior world.
For this process p-substrate is the required basic substrate, but forks and modifications will be very well possible based on a Graphene substrate or alike, still under the LSPL.
The decision to use an n-well approach is based largely on the compatibility with the existing nMOS process.
The starting material is a p-type, <100> oriented silicon with a doping concentration of $\approx 9\times10^{14}cm^{-3}$.
\label{process_overview}
\newpage
\subsection{Shallow trench isolation}\label{sti}
Shallow trench isolation (STI), also known as box isolation technique, is an integrated circuit feature which prevents electric current leakage between adjacent semiconductor device components.
STI is generally used on CMOS process technology nodes of 250 nanometers and smaller.
Older CMOS technologies and non-MOS technologies commonly use isolation based on LOCOS.
\footnote{Quirk, Michael \& Julian Serda (2001). Semiconductor Manufacturing Technology: Instructor's Manual Archived September 28, 2007, at the Wayback Machine., p. 25.} \\
The geometry of a substrate with STI implemented can be seen in \autoref{sti_target}.

\begin{figure}[H]
	\centering
	\begin{tikzpicture}[node distance = 3cm, auto, thick,scale=\CrossAndTopSectionBig, every node/.style={transform shape}]
		% substrate
\fill[substrate] (0,0) rectangle (20,2);
\node at (2,0.5) {Silicon substrate};
%trenches
\fill[isolationoxide] (0,0.75) rectangle (1,2);
\fill[isolationoxide] (8.5,0.75) rectangle (11.5,2);
\fill[isolationoxide] (19,0.75) rectangle (20,2);
	\end{tikzpicture}
	\begin{tikzpicture}[node distance = 3cm, auto, thick,scale=\CrossAndTopSectionBig, every node/.style={transform shape}]
		% substrate
\fill[YellowOrange] (0,0) rectangle (20,12);
% trench area
\fill[DarkGray] (0,0) rectangle (1,12);
\fill[DarkGray] (8.5,0) rectangle (11.5,12);
\fill[DarkGray] (19,0) rectangle (20,12);
\fill[DarkGray] (0,0) rectangle (20,1.25);
\fill[DarkGray] (0,7.5) rectangle (20,12);
	\end{tikzpicture}
	\caption{Shallow trench isolation target geometry}
	\label{sti_target}
\end{figure}

We choose the STI approach because we wanna scale the technology node down in the future below 250nm and wanna ensure backwards compatibility of our process.

\subsubsection{Pad oxide}
\subsubsection{Etching}
Dry etching (RIE)
\newpage
\section{N-well}\label{nwell_chapter}
In order to build CMOS on the same substrate, an N-well is required for building the complementary P-channel transistor for a n-p-channel logic circuitry as shown above in the example section.
The cross section as well as the top view of the targeted geometry are shown in \autoref{nwell_target}
\begin{figure}[H]
	\centering
	\begin{tikzpicture}[node distance = 3cm, auto, thick,scale=\CrossAndTopSectionBig, every node/.style={transform shape}]
		\input{tikz_process_steps/nwell.a.tex}
% p-well
\fill[pwell] (11.5,0.75) rectangle (19,2);
\node at (14.25,1) {P-Well};
% n-well
\fill[nwell] (1.25,0.75) rectangle (8.5,2);
\node at (5.75,1) {N-Well};
% p-well
\fill[pwell] (11.75,0.25) rectangle (18.75,2);
\node at (15.25,1) {P-Well};
	\end{tikzpicture}
	\begin{tikzpicture}[node distance = 3cm, auto, thick,scale=\CrossAndTopSectionBig, every node/.style={transform shape}]
		% substrate
\fill[YellowOrange] (0,0) rectangle (20,12);
% trench area
\fill[DarkGray] (0,0) rectangle (1,12);
\fill[DarkGray] (8.5,0) rectangle (11.5,12);
\fill[DarkGray] (19,0) rectangle (20,12);
\fill[DarkGray] (0,0) rectangle (20,1.25);
\fill[DarkGray] (0,7.5) rectangle (20,12);
\fill[nwell] (1.25,1) rectangle (8.25,7.25);
	\end{tikzpicture}
	\caption{N-well target geometry}
	\label{nwell_target}
\end{figure}

The N-well will serve us as an island of N-doped substrate within the P-doped basis substrate.

The dopant dose will be $2.5\times10^{12}cm^{-2}$ as calculated in the documentation of the process design leading to these steps\footnote{\url{https://github.com/leviathanch/libresiliconprocess/raw/master/process_design/process_design.pdf}}.

\begin{figure}[H]
	\centering
	\begin{tikzpicture}[node distance =1cm, auto, thick,scale=\VLSILayout, every node/.style={transform shape}]
		\input{tikz_process_steps/sti.layout.tex}
\fill[pwell,opacity=\OpacityLayout] (11.5,0.75) rectangle (19,7.5);
\fill[nwell,opacity=\OpacityLayout] (1,0.75) rectangle (8.5,7.5);
	\end{tikzpicture}
	\caption{N-Well layout}
	\label{nwell_layout}
\end{figure}

In \autoref{nwell_layout} the layout of the n-well region on top of the active area region can be seen.

The n-well is being fit into the active area. It should even be a little bit bigger than the active area, because of possible alignment offsets

\newpage

\subsection{Mask dioxide layer}
In order to selectively inject charge carrying atoms into the crystalline structure a protective dioxide ($SiO_2$) layer needs to be grown on top of a p-type substrate.
\begin{figure}[H]
	\centering
	\begin{tikzpicture}[node distance = 3cm, auto, thick,scale=\CrossSectionOnly, every node/.style={transform shape}]
		\input{tikz_process_steps/pwell.mask_dioxide_layer.a.tex}
% boron
\fill[pwell] (11.75,1.5) rectangle (18.75,2);
	\end{tikzpicture}
	\drawStepArrow{}
	\begin{tikzpicture}[node distance = 3cm, auto, thick,scale=\CrossSectionOnly, every node/.style={transform shape}]
		\input{tikz_process_steps/nwell.mask_dioxide_layer.b.tex}
	\end{tikzpicture}
	\caption{Dioxide layer growth}
\end{figure}

With an energy of 100keV for the implantation performed in \autoref{nwell_implant_step}, the projected range of the dopants within the oxide will be 100nm (130nm tops) \footnote{\url{http://cleanroom.byu.edu/rangestraggle}}.
This means being on the safe side and having 300nm as the thickness is a good approach.

In order to grow the 300nm thick oxide layer, the wafer is being oxidized for around 25 minutes at 1050\degree C using wet oxidation which results in a dioxide layer of around 300nm in thickness\footnote{\url{http://cleanroom.byu.edu/OxideTimeCalc}}.

\subsection{Patterning}

The resist is being deposited using spray coating because the uneven nature of the oxide layer.
After that the wafer is being soft baked depending on the baking time and temperature for the specific resist.
The layout for being exposed onto the resist is being extracted from the "nwell" layer within the GDS2 file onto a \textbf{bright field} mask.
The requirement is a \textbf{negative} tone resist.

\begin{figure}[H]
	\centering
	\begin{tikzpicture}[node distance = 3cm, auto, thick,scale=\CrossAndTopSection, every node/.style={transform shape}]
		% substrate
\fill[substrate] (0,0) rectangle (20,2);
\node at (2,0.5) {Silicon substrate};
%trenches
\fill[isolationoxide] (0,0.75) rectangle (1,2);
\fill[isolationoxide] (8.5,0.75) rectangle (11.5,2);
\fill[isolationoxide] (19,0.75) rectangle (20,2);

	\end{tikzpicture}
	\begin{tikzpicture}[node distance = 3cm, auto, thick,scale=\CrossAndTopSection, every node/.style={transform shape}]
		\input{tikz_process_steps/nwell.patterning.at.tex}
	\end{tikzpicture}
	\drawStepArrow{Mask: nwell}
	\begin{tikzpicture}[node distance = 3cm, auto, thick,scale=\CrossAndTopSection, every node/.style={transform shape}]
		% resist
\fill[resist] (0,2.6) rectangle (1,5.0);
\fill[resist] (8.5,2.6) rectangle (20,5.0);

\input{tikz_process_steps/nwell.mask_dioxide_layer.b.tex}
	\end{tikzpicture}
	\begin{tikzpicture}[node distance = 3cm, auto, thick,scale=\CrossAndTopSection, every node/.style={transform shape}]
		\input{tikz_process_steps/nwell.patterning.bt.tex}
	\end{tikzpicture}
	\caption{Cross/top view of n-well layout on resist}
\end{figure}

The thickness of the resist layer and the baking duration will variate depending on the specific equipment for which this process will be implemented with.
Also after the exposure and development, the hard baking shouldn't be forgotten!

\newpage

\subsection{Etching}
We now need to open a window in the dioxide layer, through which we will inject carrier atoms into the silicon crystal structure.
\begin{figure}[H]
	\centering
	\begin{tikzpicture}[node distance = 3cm, auto, thick,scale=\CrossAndTopSection, every node/.style={transform shape}]
		% resist
\fill[resist] (0,2.6) rectangle (1,5.0);
\fill[resist] (8.5,2.6) rectangle (20,5.0);

\input{tikz_process_steps/nwell.mask_dioxide_layer.b.tex}
	\end{tikzpicture}
	\begin{tikzpicture}[node distance = 3cm, auto, thick,scale=\CrossAndTopSection, every node/.style={transform shape}]
		% resist
\fill[resist] (0,0) rectangle (20,12);
% substrate
\fill[isolationoxide]  (1,0.75) rectangle (8.5,7.5);
	\end{tikzpicture}
	\drawStepArrow{}
	\begin{tikzpicture}[node distance = 3cm, auto, thick,scale=\CrossAndTopSection, every node/.style={transform shape}]
		\input{tikz_process_steps/nwell.etching.b.tex}
	\end{tikzpicture}
	\begin{tikzpicture}[node distance = 3cm, auto, thick,scale=\CrossAndTopSection, every node/.style={transform shape}]
		\input{tikz_process_steps/nwell.etching.bt.tex}
	\end{tikzpicture}
	\caption{Cross/top view of n-well oxide window}
\end{figure}

Since the silicon dioxide layer is 300nm thick and we wanna reach the silicon below we can use wet etching as described in the chemistry chapter.\\

\textbf{Possible approaches}:
\begin{itemize}
	\item \textbf{"AOE Etcher (DRY-AOE)" from HKUST} \\
	We can use anisotropic plasma etching for sharper borders.
	\item \textbf{Chemical solution} \\
	We can use buffered hydrofluoric acid (BOE (1:6)) at room temperature for a little bit over 3 minutes in order to get through the 300nm of oxide.\\
	Too long over 3 minutes might cause under-etch however!
\end{itemize}

\subsection{Cleaning}
In order to avoid contamination of the machines we need to make sure all the resist has been stripped off from the wafer.
\begin{figure}[H]
	\centering
	\begin{tikzpicture}[node distance = 3cm, auto, thick,scale=\CrossSectionOnly, every node/.style={transform shape}]
		\input{tikz_process_steps/nwell.etching.b.tex}
	\end{tikzpicture}
	\drawStepArrow{}
	\begin{tikzpicture}[node distance = 3cm, auto, thick,scale=\CrossSectionOnly, every node/.style={transform shape}]
		% oxide
\fill[isolationoxide] (0,1.25) rectangle (1,2.75);
\fill[isolationoxide] (8.5,1.25) rectangle (20,2.75);

% oxide hill 1
\filldraw[line width=0, isolationoxide] (0.5,2.75) -- (1.0,2.75) -- (1.0,3.0);
\filldraw[line width=0, isolationoxide] (8.75,2.75) -- (8.5,2.75) -- (8.5,3.0);

% oxide hill 2
\fill[isolationoxide] (11.75,2.75) rectangle (18.75,3.5);
\filldraw[line width=0, isolationoxide] (11.0,2.75) -- (11.75,2.75) -- (11.75,3.5);
\filldraw[line width=0, isolationoxide] (18.75,2.75) -- (18.75,3.5) -- (19.5,2.75);

\input{tikz_process_steps/pwell.oxide_mask_removal.b.tex}

	\end{tikzpicture}
	\caption{Resist removal}
\end{figure}
Please just use the solvent for the specific resist.

\newpage

\subsection{Implantation/Doping}\label{nwell_implant_step}
We now need to inject the carriers into the upper level of the n-channel area so that we can later on drive them into the crystal during the drive-in step.

\begin{figure}[H]
	\centering
	\begin{minipage}{0.5\textwidth}
	\centering
	\begin{tikzpicture}[node distance = 3cm, auto, thick,scale=\CrossSectionOnly, every node/.style={transform shape}]
		% resist
\fill[resist] (0,2.6) rectangle (1,5.0);
\fill[resist] (8.5,2.6) rectangle (20,5.0);

\input{tikz_process_steps/nwell.mask_dioxide_layer.b.tex}

\forloop{ct}{0}{\value{ct} < 21}
{
	\draw [->] (\value{ct},6) -- (\value{ct},5);
	\node at (\value{ct},6.2) {P$^{31}$};
}

	\end{tikzpicture}
	\drawStepArrow{Boron implant}
	\begin{tikzpicture}[node distance = 3cm, auto, thick,scale=\CrossSectionOnly, every node/.style={transform shape}]
		% oxide
\fill[isolationoxide] (0,1.25) rectangle (1,2.75);
\fill[isolationoxide] (8.5,1.25) rectangle (20,2.75);

% oxide hill 1
\filldraw[line width=0, isolationoxide] (0.5,2.75) -- (1.0,2.75) -- (1.0,3.0);
\filldraw[line width=0, isolationoxide] (8.75,2.75) -- (8.5,2.75) -- (8.5,3.0);

% oxide hill 2
\fill[isolationoxide] (11.75,2.75) rectangle (18.75,3.5);
\filldraw[line width=0, isolationoxide] (11.0,2.75) -- (11.75,2.75) -- (11.75,3.5);
\filldraw[line width=0, isolationoxide] (18.75,2.75) -- (18.75,3.5) -- (19.5,2.75);

\input{tikz_process_steps/nwell.mask_dioxide_layer.a.tex}


% phosphorus
\fill[nwell] (1.25,1.8) rectangle (8.25,2);
\node at (4.75,1.9) {n+};

	\end{tikzpicture}
	\drawStepArrow{Annealing}
	\begin{tikzpicture}[node distance = 3cm, auto, thick,scale=\CrossSectionOnly, every node/.style={transform shape}]
		% oxide
\fill[isolationoxide] (0,1.25) rectangle (1,2.75);
\fill[isolationoxide] (8.5,1.25) rectangle (20,2.75);

% oxide hill 1
\filldraw[line width=0, isolationoxide] (0.5,2.75) -- (1.0,2.75) -- (1.0,3.0);
\filldraw[line width=0, isolationoxide] (8.75,2.75) -- (8.5,2.75) -- (8.5,3.0);

% oxide hill 2
\fill[isolationoxide] (11.75,2.75) rectangle (18.75,3.5);
\filldraw[line width=0, isolationoxide] (11.0,2.75) -- (11.75,2.75) -- (11.75,3.5);
\filldraw[line width=0, isolationoxide] (18.75,2.75) -- (18.75,3.5) -- (19.5,2.75);

\input{tikz_process_steps/nwell.mask_dioxide_layer.a.tex}

\fill[nwell] (1.25,1.5) rectangle (8.25,2.0);
\fill[pwell] (11.75,1.0) rectangle (18.75,2.0);
	\end{tikzpicture} \\
	\textbf{Implantation approach}
	\end{minipage}\begin{minipage}{0.5\textwidth}
	\centering
	\begin{tikzpicture}[node distance = 3cm, auto, thick,scale=\CrossSectionOnly, every node/.style={transform shape}]
		\fill[nwell] (0,0) rectangle (20,5);

% oxide
\fill[isolationoxide] (0,1.25) rectangle (1,2.75);
\fill[isolationoxide] (8.5,1.25) rectangle (20,2.75);

% oxide hill 1
\filldraw[line width=0, isolationoxide] (0.5,2.75) -- (1.0,2.75) -- (1.0,3.0);
\filldraw[line width=0, isolationoxide] (8.75,2.75) -- (8.5,2.75) -- (8.5,3.0);

% oxide hill 2
\fill[isolationoxide] (11.75,2.75) rectangle (18.75,3.5);
\filldraw[line width=0, isolationoxide] (11.0,2.75) -- (11.75,2.75) -- (11.75,3.5);
\filldraw[line width=0, isolationoxide] (18.75,2.75) -- (18.75,3.5) -- (19.5,2.75);

\input{tikz_process_steps/nwell.mask_dioxide_layer.a.tex}

	\end{tikzpicture}
	\drawStepArrow{Constant source diffusion}
	\begin{tikzpicture}[node distance = 3cm, auto, thick,scale=\CrossSectionOnly, every node/.style={transform shape}]
		\input{tikz_process_steps/nwell.doping.b.tex}
	\end{tikzpicture}
	\drawStepArrow{Source removal}
	\begin{tikzpicture}[node distance = 3cm, auto, thick,scale=\CrossSectionOnly, every node/.style={transform shape}]
		% oxide
\fill[isolationoxide] (0,1.25) rectangle (1,2.75);
\fill[isolationoxide] (8.5,1.25) rectangle (20,2.75);

% oxide hill 1
\filldraw[line width=0, isolationoxide] (0.5,2.75) -- (1.0,2.75) -- (1.0,3.0);
\filldraw[line width=0, isolationoxide] (8.75,2.75) -- (8.5,2.75) -- (8.5,3.0);

% oxide hill 2
\fill[isolationoxide] (11.75,2.75) rectangle (18.75,3.5);
\filldraw[line width=0, isolationoxide] (11.0,2.75) -- (11.75,2.75) -- (11.75,3.5);
\filldraw[line width=0, isolationoxide] (18.75,2.75) -- (18.75,3.5) -- (19.5,2.75);

\input{tikz_process_steps/nwell.mask_dioxide_layer.a.tex}


\shade[upper left = nwell, upper right = nwell, lower right = substrate, lower left = substrate,] (1.25,1.5) rectangle (8.25,2.0);
\shade[upper left = pwell, upper right = pwell, lower right = substrate, lower left = substrate,] (11.75,1.0) rectangle (18.75,2.0);
	\end{tikzpicture} \\
	\textbf{Diffusion approach}
	\end{minipage}
	\caption{Doping process}
\end{figure}


\textbf{Possible approaches}:
\begin{itemize}
	\item \textbf{"CF-3000 Implanter (IMP-3000)" from HKUST} \\
	At HKUST we have an implanter which gives us better control over the initial surface concentration. \\
	These steps are needed to arrive with the desired geometry:
	\begin{enumerate}
		\item Preparing by default cleaning
		\item The N-well is implanted with a Phosphorus ($P^{31}$) dose of $2.5\times10^{12}cm^{-2}$ at an energy of 100 keV.
		\item The N-well is annealed for 30 minutes at 1050\degreesC in $N_2$ environment (DIF-A1)\\
		After that the P-well will be around 2\um deep and the N-Well around 1\um deep
	\end{enumerate}
	\item \textbf{Constant source diffusion} \\
	We can add a layer of Phosphorus solution and diffusing in order to have an initial concentration in order to reach the desired concentration later by main diffusion.
		\begin{enumerate}
		\item A constant source is added (gas or liquid)
		\item The source dopant is driven in for 10 minutes at 1050\degreesC
		\item The dopant source is removed by stopping the gas flow or cleaning the surface
	\end{enumerate}
\end{itemize}

\newpage

\subsection{Oxide for drive-in}

Now we need to cover the now doped and annealed area with an oxide layer in order to seal it off and prevent the dopants from diffusing out of the crystal during the drive-in phase.

\begin{figure}[H]
	\centering
	\begin{tikzpicture}[node distance = 3cm, auto, thick,scale=\CrossSectionOnly, every node/.style={transform shape}]
		\input{tikz_process_steps/nwell.drive_in_oxide.a.tex}
	\end{tikzpicture}
	\drawStepArrow{}
	\begin{tikzpicture}[node distance = 3cm, auto, thick,scale=\CrossSectionOnly, every node/.style={transform shape}]
		\fill[isolationoxide] (0,0) rectangle (20,2.25);
\input{tikz_process_steps/nwell.cleaning.b.tex}

% phosphorus
\fill[nwell] (1.25,1.8) rectangle (8.25,2);
\node at (4.75,1.9) {n+};

	\end{tikzpicture}
	\caption{Oxide growth}
\end{figure}

The wafer is being oxidized for around 5 minutes at 1050\degree C in a wet environment in order to achieve a cover silicon layer of around 90nm thickness.

\subsection{Drive-in}
In order to drive the carrier atoms deeper into the crystalline structure the wafer needs to be driven in after predeposition.

\begin{figure}[H]
	\centering
	\begin{tikzpicture}[node distance = 3cm, auto, thick,scale=\CrossSectionOnly, every node/.style={transform shape}]
		\fill[isolationoxide] (0,0) rectangle (20,2.25);
\input{tikz_process_steps/nwell.implantation.b.tex}
	\end{tikzpicture}
	\drawStepArrow{}
	\begin{tikzpicture}[node distance = 3cm, auto, thick,scale=\CrossSectionOnly, every node/.style={transform shape}]
		\fill[isolationoxide] (0,0) rectangle (20,2.25);
\input{tikz_process_steps/nwell.cleaning.b.tex}

% phosphorus
\fill[nwell] (1.25,1.8) rectangle (8.25,2);
\node at (4.75,1.9) {n+};


% n-well
\fill[nwell] (1.25,0.75) rectangle (8.25,2);
\node at (4.75,1) {N-Well};
	\end{tikzpicture}
	\caption{Drive-in process}
\end{figure}

In this step the wafer is  driven-in for 90 minutes at 1050\degree C in an inert ambient\footnote{\url{http://cleanroom.byu.edu/DopConCalc}}

\subsection{Oxide mask removal}
Now we want to remove the silicon mask from the wafer and clean it for another clean oxide mask layer.

\begin{figure}[H]
	\centering
	\begin{tikzpicture}[node distance = 3cm, auto, thick,scale=\CrossSectionOnly, every node/.style={transform shape}]
		\input{tikz_process_steps/nwell.cleaning.b.tex}
\fill[nwell] (1.25,1.5) rectangle (8.25,2.0);
\fill[pwell] (11.75,1.0) rectangle (18.75,2.0);
	\end{tikzpicture}
	\drawStepArrow{}
	\begin{tikzpicture}[node distance = 3cm, auto, thick,scale=\CrossSectionOnly, every node/.style={transform shape}]
		\input{tikz_process_steps/nwell.a.tex}
% p-well
\fill[pwell] (11.5,0.75) rectangle (19,2);
\node at (14.25,1) {P-Well};
% n-well
\shade[upper left = nwell, upper right = nwell, lower right = substrate, lower left = substrate,] (1.25,0.75) rectangle (8.5,2);
\node at (5.75,1) {N-Well};
% p-well
\shade[upper left = pwell, upper right = pwell, lower right = substrate, lower left = substrate,] (11.75,0.25) rectangle (18.75,2);
\node at (15.25,1) {P-Well};
	\end{tikzpicture}
	\caption{Oxide removal}
\end{figure}

We use buffered hydrofluoric acid (BOE (1:6)) at room temperature for 3 minutes in order to remove the 300nm of oxide layer.
\newpage
`\section{P-well}\label{pwell_chapter}
In order to build CMOS on the same substrate, a P-well is required for building the complementary N-channel transistor for a n-p-channel logic circuitry.
The cross section as well as the top view of the targeted geometry are shown in \autoref{nwell_target}
\begin{figure}[H]
	\centering
	\begin{tikzpicture}[node distance = 3cm, auto, thick,scale=\CrossAndTopSectionBig, every node/.style={transform shape}]
		\input{tikz_process_steps/pwell.a.tex}
% n-well
\fill[nwell] (1.25,0.75) rectangle (8.5,2);
\node at (5.75,1) {N-Well};
% p-well
\fill[pwell] (11.75,0.25) rectangle (18.75,2);
\node at (15.25,1) {P-Well};
% p-well
\fill[pwell] (11.5,0.75) rectangle (19,2);
\node at (14.25,1) {P-Well};
	\end{tikzpicture}
	\begin{tikzpicture}[node distance = 3cm, auto, thick,scale=\CrossAndTopSectionBig, every node/.style={transform shape}]
		\input{tikz_process_steps/sti.b.tex}
\fill[nwell] (1.25,1) rectangle (8.25,7.25);
\fill[pwell] (11.5,1.25) rectangle (19,7.5);
	\end{tikzpicture}
	\caption{P-well target geometry}
	\label{pwell_target}
\end{figure}
The P-well will serve us as an island of higher p-doped substrate within the slightly p-doped basis substrate.

The dopant dose will be $2.5\times10^{12}cm^{-2}$ as calculated in the documentation of the process design leading to these steps\footnote{\url{https://github.com/leviathanch/libresiliconprocess/raw/master/process_design/process_design.pdf}}.

\begin{figure}[H]
	\centering
	\begin{tikzpicture}[node distance =1cm, auto, thick,scale=\VLSILayout, every node/.style={transform shape}]
		\fill[Goldenrod,opacity=0.2] (0.75,0.5) rectangle (8.75,7.75);
\fill[Goldenrod,opacity=0.2] (11.25,0.5) rectangle (19.25,7.75);

\draw[dotted] (20.5,0.5) rectangle (25,5.5);

\node at (22.25,5) {\textbf{Layers}};

\fill[Goldenrod,opacity=0.2] (21,1) rectangle (21.5,1.5);
\node at (22.25,1.25) {active};

\fill[orange,opacity=0.2] (21,1.5) rectangle (21.5,2);
\node at (22.25,1.75) {nwell};

\fill[blue,opacity=0.2] (21,2) rectangle (21.5,2.5);
\node at (22.25,2.25) {nimplant};

\fill[red,opacity=0.2] (21,2.5) rectangle (21.5,3);
\node at (22.25,2.75) {pimplant};

\fill[Emerald,opacity=0.2] (21,3) rectangle (21.5,3.5);
\node at (22.25,3.25) {gate};

\fill[Fuchsia,opacity=0.2] (21,3.5) rectangle (21.5,4);
\node at (22.25,3.75) {metal1};

\fill[DarkOrchid,opacity=0.2] (21,4) rectangle (21.5,4.5);
\node at (22.25,4.25) {via1};
\fill[pwell,opacity=\OpacityLayout] (11.5,0.75) rectangle (19,7.5);
	\end{tikzpicture}
	\caption{P-Well layout}
	\label{pwell_layout}
\end{figure}

In \autoref{pwell_layout} the layout of the P-well region on top of the active area region can be seen.

The p-well is being fit into the active area.

It should even be a little bit bigger than the active area, because of possible alignment offsets

\newpage

\subsection{Mask dioxide layer}
In order to selectively inject charge carrying atoms into the crystalline structure a protective dioxide ($SiO_2$) layer needs to be grown on top of a p-type substrate.
\begin{figure}[H]
	\centering
	\begin{tikzpicture}[node distance = 3cm, auto, thick,scale=\CrossSectionOnly, every node/.style={transform shape}]
		\input{tikz_process_steps/pwell.mask_dioxide_layer.a.tex}
	\end{tikzpicture} \\
	\includegraphics[scale=0.01]{down_arrow.png} \\
	\begin{tikzpicture}[node distance = 3cm, auto, thick,scale=\CrossSectionOnly, every node/.style={transform shape}]
		\input{tikz_process_steps/pwell.mask_dioxide_layer.b.tex}
	\end{tikzpicture}
	\caption{Dioxide layer growth}
\end{figure}
With an energy of 100keV for the implantation performed in \autoref{pwell_implant_step}, the projected range of the dopants within the oxide will be 310nm (380nm tops) \footnote{\url{http://cleanroom.byu.edu/rangestraggle}}.
This means being on the safe side and having 500nm as the thickness is a good approach.

In order to grow the 500nm thick oxide layer, the wafer is being oxidized for around 56 minutes at 1050\degree C using wet oxidation which results in a dioxide layer of around 500nm in thickness\footnote{\url{http://cleanroom.byu.edu/OxideTimeCalc}}.

\subsection{Patterning}
The resist is being deposited spray or spin coating (spray coating is better because of the uneven surface!) and then soft baked depending on the baking time for the specific resist.
The layout for being exposed onto the resist is being extracted from the "pwell" layer within the GDS2 file.
\begin{figure}[H]
	\centering
	\begin{tikzpicture}[node distance = 3cm, auto, thick,scale=\CrossAndTopSection, every node/.style={transform shape}]
		% substrate
\fill[substrate] (0,0) rectangle (20,1.25);
\node at (2,0.5) {Silicon substrate};
\fill[substrate] (0.25,1.25) rectangle (19.75,2);


	\end{tikzpicture}
	\begin{tikzpicture}[node distance = 3cm, auto, thick,scale=\CrossAndTopSection, every node/.style={transform shape}]
		\input{tikz_process_steps/pwell.patterning.at.tex}
	\end{tikzpicture} \\
	\includegraphics[scale=0.01]{down_arrow.png} \\
	\begin{tikzpicture}[node distance = 3cm, auto, thick,scale=\CrossAndTopSection, every node/.style={transform shape}]
		% resist
\fill[resist] (0.25,2.0) rectangle (11.5,5.0);
\fill[resist] (19,2.0) rectangle (19.75,5.0);

\input{tikz_process_steps/basic.a.tex}


	\end{tikzpicture}
	\begin{tikzpicture}[node distance = 3cm, auto, thick,scale=\CrossAndTopSection, every node/.style={transform shape}]
		\input{tikz_process_steps/pwell.patterning.bt.tex}
	\end{tikzpicture}
	\caption{Cross/top view of P-well layout on resist}
\end{figure}
The thickness of the resist layer and the baking duration will variate depending on the specific equipment for which this process will be implemented with.
Also after the exposure and development, the hard baking shouldn't be forgotten!

\newpage

\subsection{Etching}
We now need to open a window in the dioxide layer, through which we will inject carrier atoms into the silicon crystal structure.
\begin{figure}[H]
	\centering
	\begin{tikzpicture}[node distance = 3cm, auto, thick,scale=\CrossAndTopSection, every node/.style={transform shape}]
		% resist
\fill[resist] (0.25,2.0) rectangle (11.5,5.0);
\fill[resist] (19,2.0) rectangle (19.75,5.0);

\input{tikz_process_steps/pwell.patterning.a.tex}

	\end{tikzpicture}
	\begin{tikzpicture}[node distance = 3cm, auto, thick,scale=\CrossAndTopSection, every node/.style={transform shape}]
		% resist
\fill[resist] (0,0) rectangle (20,12);
% substrate
\fill[isolationoxide] (11.5,1.5) rectangle (19,7.25);
	\end{tikzpicture} \\
	\includegraphics[scale=0.01]{down_arrow.png} \\
	\begin{tikzpicture}[node distance = 3cm, auto, thick,scale=\CrossAndTopSection, every node/.style={transform shape}]
		% resist
\fill[resist] (0,2.6) rectangle (11.5,5.0);
\fill[resist] (19,2.6) rectangle (20,5.0);

% oxide
\fill[isolationoxide] (0,1.25) rectangle (11.5,2.75);
\fill[isolationoxide] (19,1.25) rectangle (20,2.75);

% oxide hill 1
\fill[isolationoxide] (1.25,2.75) rectangle (8.25,3.5);
\filldraw[line width=0, isolationoxide] (0.5,2.75) -- (1.25,2.75) -- (1.25,3.5);
\filldraw[line width=0, isolationoxide] (8.25,2.75) -- (8.25,3.5) -- (9.0,2.75);

% oxide hill 2
\filldraw[line width=0, isolationoxide] (11.25,2.75) -- (11.5,2.75) -- (11.5,3.0);
\filldraw[line width=0, isolationoxide] (19.0,3.0)  -- (19.0,2.75) -- (19.25,2.75);

\node at (2,2.1) {SiO2};

\input{tikz_process_steps/pwell.mask_dioxide_layer.a.tex}


	\end{tikzpicture}
	\begin{tikzpicture}[node distance = 3cm, auto, thick,scale=\CrossAndTopSection, every node/.style={transform shape}]
		% resist
\fill[resist] (0,0) rectangle (20,12);
% substrate
\fill[substrate] (11.5,1.5) rectangle (19,7.25);
	\end{tikzpicture}
	\caption{Cross/top view of P-well oxide window}
\end{figure}

There are multiple possible approaches to etch through these 500nm of oxide.

\textbf{Possible approaches}:
\begin{itemize}
	\item \textbf{"DRIE Etcher \#1" from HKUST} \\
	We can use anisotropic plasma etching for sharper borders.
	\item \textbf{Chemical solution} \\
	We can use buffered hydrofluoric acid (BOE (1:6)) at room temperature for a little bit over 4 minutes in order to get through the 500nm of oxide.\\
	Too long over 4 minutes might cause under-etch however!
\end{itemize}

\subsection{Cleaning}
In order to avoid contamination of the machines we need to make sure all the resist has been stripped off from the wafer.
\begin{figure}[H]
	\centering
	\begin{tikzpicture}[node distance = 3cm, auto, thick,scale=\CrossSectionOnly, every node/.style={transform shape}]
		% oxide
\fill[isolationoxide] (0,1.25) rectangle (11.5,2.75);
\fill[isolationoxide] (19,1.25) rectangle (20,2.75);

% oxide hill 1
\fill[isolationoxide] (1.25,2.75) rectangle (8.25,3.5);
\filldraw[line width=0, isolationoxide] (0.5,2.75) -- (1.25,2.75) -- (1.25,3.5);
\filldraw[line width=0, isolationoxide] (8.25,2.75) -- (8.25,3.5) -- (9.0,2.75);

% oxide hill 2
\filldraw[line width=0, isolationoxide] (11.25,2.75) -- (11.5,2.75) -- (11.5,3.0);
\filldraw[line width=0, isolationoxide] (19.0,3.0)  -- (19.0,2.75) -- (19.25,2.75);

\node at (2,2.1) {SiO2};

\input{tikz_process_steps/pwell.mask_dioxide_layer.a.tex}

% boron
\fill[pwell] (11.75,1.8) rectangle (18.75,2);

	\end{tikzpicture} \\
	\includegraphics[scale=0.01]{down_arrow.png} \\
	\begin{tikzpicture}[node distance = 3cm, auto, thick,scale=\CrossSectionOnly, every node/.style={transform shape}]
		% substrate
\fill[substrate] (0,0) rectangle (20,1.25);
\node at (2,0.5) {Silicon substrate};
\fill[substrate] (0.25,1.25) rectangle (19.75,2);

% boron
\shade[upper left = pwell, upper right = pwell, lower right = substrate, lower left = substrate,] (11.5,1.5) rectangle (19.0,2);

	\end{tikzpicture}
	\caption{Resist removal}
\end{figure}
Please just use the solvent for the specific resist.

\newpage

\subsection{Implantation/Predeposition}\label{pwell_implant_step}
We now need to inject the carriers into the upper level of the n-channel area so that we can later on drive them into the crystal during the drive-in step.
\begin{figure}[H]
	\centering
	\begin{tikzpicture}[node distance = 3cm, auto, thick,scale=\CrossSectionOnly, every node/.style={transform shape}]
		\input{tikz_process_steps/basic.a.tex}
% boron
\shade[upper left = pwell, upper right = pwell, lower right = substrate, lower left = substrate,] (11.5,1.5) rectangle (19.0,2);


\forloop{ct}{0}{\value{ct} < 21}
{
	\draw [->] (\value{ct},5) -- (\value{ct},4);
	\node at (\value{ct},5.2) {B$^{11}$};
}
	\end{tikzpicture} \\
	\includegraphics[scale=0.01]{down_arrow.png} \\
	\begin{tikzpicture}[node distance = 3cm, auto, thick,scale=\CrossSectionOnly, every node/.style={transform shape}]
		% oxide
\fill[isolationoxide] (0,1.25) rectangle (11.5,2.75);
\fill[isolationoxide] (19,1.25) rectangle (20,2.75);

% oxide hill 1
\fill[isolationoxide] (1.25,2.75) rectangle (8.25,3.5);
\filldraw[line width=0, isolationoxide] (0.5,2.75) -- (1.25,2.75) -- (1.25,3.5);
\filldraw[line width=0, isolationoxide] (8.25,2.75) -- (8.25,3.5) -- (9.0,2.75);

% oxide hill 2
\filldraw[line width=0, isolationoxide] (11.25,2.75) -- (11.5,2.75) -- (11.5,3.0);
\filldraw[line width=0, isolationoxide] (19.0,3.0)  -- (19.0,2.75) -- (19.25,2.75);

\node at (2,2.1) {SiO2};

\input{tikz_process_steps/pwell.mask_dioxide_layer.a.tex}

% boron
\fill[pwell] (11.75,1.8) rectangle (18.75,2);
	\end{tikzpicture}
	\caption{Doping process}
\end{figure}

\textbf{Possible approaches}:
\begin{itemize}
	\item \textbf{"CF-3000 Implanter (IMP-3000)" from HKUST} \\
	The P-well is implanted with a Boron ($B^{11}$) dose of $2.5\times10^{12}cm^{-2}$ at an energy of 100 KeV.\\
	The P-well is then annealed.
	\item \textbf{Constant source diffusion} \\
	We can add a layer of Boron solution and diffusing in order to have an initial concentration in order to reach the desired concentration later by main diffusion.
\end{itemize}

\subsection{Oxide mask removal}
We want to remove the silicon mask from the wafer so that the P-well becomes accessible for the further process steps but we don't want to etch "way too much" of the trench material.
\begin{figure}[H]
	\centering
	\begin{tikzpicture}[node distance = 3cm, auto, thick,scale=\CrossSectionOnly, every node/.style={transform shape}]
		\input{tikz_process_steps/pwell.implantation.c.tex}
	\end{tikzpicture} \\
	\includegraphics[scale=0.01]{down_arrow.png} \\
	\begin{tikzpicture}[node distance = 3cm, auto, thick,scale=\CrossSectionOnly, every node/.style={transform shape}]
		\input{tikz_process_steps/pwell.mask_dioxide_layer.a.tex}
% boron
\fill[pwell] (11.75,1.5) rectangle (18.75,2);
	\end{tikzpicture}
	\caption{Oxide removal}
\end{figure}
We use buffered hydrofluoric acid (BOE (1:6)) at room temperature for a little bit over 4 minutes (roughly 10 seconds more or so) in order to remove the 500nm of oxide layer.
\newpage
\section{n+ Implant}\label{nimplant}
For the bulk of the PMOS transistors and for the source and drain of the NMOS transistors highly doped  n+ areas are required.
In this step we're going to build these.

\begin{figure}[H]
	\centering
	\begin{tikzpicture}[node distance = 3cm, auto, thick,scale=\CrossAndTopSectionBig, every node/.style={transform shape}]
		% substrate
\fill[YellowOrange] (0,0) rectangle (20,2);
\node at (2,0.5) {Si (p-type)};
% n-well
\fill[Goldenrod] (1.25,0.75) rectangle (8.25,2);
\node at (5.75,1) {N-Well};
% body
\fill[ProcessBlue] (1.5,1) rectangle (3,2);
\node at (2,1.5) {n+};
% source
\fill[RedOrange] (3.5,1) rectangle (5,2);
\node at (4,1.5) {p+};
% drain
\fill[RedOrange] (6.5,1) rectangle (8,2);
\node at (7,1.5) {p+};
%% gate:
% gate oxide
\fill[LightGray] (4.8,2) rectangle (6.7,2.1);
% gate poly
\fill[BrickRed] (4.8,2.1) rectangle (6.7,2.2);

%field oxides:
\fill[DarkGray] (0,2) rectangle (1,4);
\fill[DarkGray] (8.5,2) rectangle (11.5,4);
\fill[DarkGray] (19,2) rectangle (20,4);

\fill[RedOrange] (0,1.5) rectangle (1,2);
\fill[RedOrange] (8.5,1.5) rectangle (11.5,2);
\fill[RedOrange] (19,1.5) rectangle (20,2);

\node at (0.5,1.75) {p+};
\node at (9.5,1.75) {p+};
\node at (19.5,1.75) {p+};

%%% nmos:
% body
\fill[RedOrange] (17,1) rectangle (18.5,2);
\node at (18,1.5) {p+};
% source
\fill[ProcessBlue] (15,1) rectangle (16.5,2);
\node at (16,1.5) {n+};
% drain
\fill[ProcessBlue] (12,1) rectangle (13.5,2);
\node at (13,1.5) {n+};

%% gate:
% gate oxide
\fill[LightGray] (13.3,2) rectangle (15.2,2.1);
% gate poly
\fill[BrickRed] (13.3,2.1) rectangle (15.2,2.2);
	\end{tikzpicture}
	\begin{tikzpicture}[node distance = 3cm, auto, thick,scale=\CrossAndTopSectionBig, every node/.style={transform shape}]
		\fill[isolationoxide] (0,0) rectangle (20,10);

% n-well
\fill[nwell] (1,1.25) rectangle (8.5,7.5);

% p-well
\fill[pwell] (11.5,1.25) rectangle (19,7.5);

% gate metal
\fill[gatemetal] (5,0) rectangle (6.5,9);
\fill[gatemetal] (13.5,0) rectangle (15,9);
\fill[gatemetal] (5,8) rectangle (15,10);

% n+
\fill[nimplant] (1.5,2) rectangle (3,6.5);
\fill[nimplant] (12,2) rectangle (13.5,6.5);
\fill[nimplant] (15,2) rectangle (16.5,6.5);

	\end{tikzpicture}
	\caption{N+ implant geometry target}
\end{figure}

The tricky thing here is to have a reasonable implant depth but not too deep because the deeper the junction, the higher the junction capacity which in turn limits the switching performance of the CMOS circuitry.

\begin{figure}[H]
	\centering
	\begin{tikzpicture}[node distance =1cm, auto, thick,scale=\VLSILayout, every node/.style={transform shape}]
		\input{tikz_process_steps/pimplant.layout.tex}
% gate metal
\fill[gatemetal,opacity=\OpacityLayout] (4.8,1.75) rectangle (6.7,8);
\fill[gatemetal,opacity=\OpacityLayout] (13.3,1.75) rectangle (15.2,8);
\fill[gatemetal,opacity=\OpacityLayout] (4.8,8) rectangle (15.2,10);


% n+
\fill[nimplant,opacity=\OpacityLayout] (1.5,2) rectangle (3,6.5);
\fill[nimplant,opacity=\OpacityLayout] (12,2) rectangle (13.5,6.5);
\fill[nimplant,opacity=\OpacityLayout] (15,2) rectangle (16.5,6.5);

	\end{tikzpicture}
	\caption{N+ layout}
	\label{nimplant_layout}
\end{figure}

An example layout of p-implants can be seen in \autoref{nimplant_layout}, the mask is being extracted from the layer "n\_plus\_select".

Also important to notice is that this example layout is just for demonstration purposes only, please have a look at the standard cell documentation for the actual layouts. 

\newpage

\subsection{Mask dioxide layer}

In order to selectively inject charge carrying atoms into the crystalline structure a protective dioxide ($SiO_2$) layer needs to be grown on top of a p-type substrate.

\begin{figure}[H]
	\centering
	\begin{tikzpicture}[node distance = 3cm, auto, thick,scale=\CrossSectionOnly, every node/.style={transform shape}]
		\input{tikz_process_steps/nimplant.oxide_growth.a.tex}
	\end{tikzpicture}
	\drawStepArrow{}
	\begin{tikzpicture}[node distance = 3cm, auto, thick,scale=\CrossSectionOnly, every node/.style={transform shape}]
		\input{tikz_process_steps/nimplant.oxide_growth.b.tex}
	\end{tikzpicture}
	\caption{Oxide layer}
\end{figure}

With an energy of 100keV for the implantation performed in \autoref{pwell_implant_step}, the projected range of the dopants within the oxide will be 310nm (380nm tops) \footnote{\url{http://cleanroom.byu.edu/rangestraggle}}.
This means being on the safe side and having 500nm as the thickness is a good approach.
In order to grow the 500nm thick oxide layer, the wafer is being oxidized for around 56 minutes at 1050\degree C using wet oxidation which results in a dioxide layer of around 500nm in thickness\footnote{\url{http://cleanroom.byu.edu/OxideTimeCalc}}.

\subsection{Pattering}
\begin{figure}[H]
	\centering
	\begin{tikzpicture}[node distance = 3cm, auto, thick,scale=\CrossAndTopSection, every node/.style={transform shape}]
		\input{tikz_process_steps/pwell.a.tex}
\fill[gateoxide] (4.8,2) rectangle (6.7,2.3);
\fill[gateoxide] (13.3,2) rectangle (15.2,2.3);
\fill[gatemetal] (4.8,2.3) rectangle (6.7,2.6);
\fill[gatemetal] (13.3,2.3) rectangle (15.2,2.6);

	\end{tikzpicture}
	\begin{tikzpicture}[node distance = 3cm, auto, thick,scale=\CrossAndTopSection, every node/.style={transform shape}]
		\input{tikz_process_steps/nimplant.patterning.at.tex}
	\end{tikzpicture}
	\drawStepArrow{}
	\begin{tikzpicture}[node distance = 3cm, auto, thick,scale=\CrossAndTopSection, every node/.style={transform shape}]
		% resist
\fill[resist] (0,2.0) rectangle (1.25,6.0);
\fill[resist] (3,2.0) rectangle (11.75,6.0);
\fill[resist] (17,2.0) rectangle (20,6.0);

\input{tikz_process_steps/nimplant.oxide_growth.b.tex}

	\end{tikzpicture}
	\begin{tikzpicture}[node distance = 3cm, auto, thick,scale=\CrossAndTopSection, every node/.style={transform shape}]
		\fill[resist] (0,0) rectangle (20,12);

% n+
\fill[isolationoxide] (1.5,2) rectangle (3,6.5);
\fill[isolationoxide] (12,2) rectangle (13.5,6.5);
\fill[isolationoxide] (15,2) rectangle (16.5,6.5);
	\end{tikzpicture}
	\caption{N+ region resist mask}
\end{figure}

\subsection{Etching}
\begin{figure}[H]
	\centering
	\begin{tikzpicture}[node distance = 3cm, auto, thick,scale=\CrossAndTopSection, every node/.style={transform shape}]
		\input{tikz_process_steps/well.a.tex}
% oxide
\fill[isolationoxide] (0,2) rectangle (20,3);
% resist
\fill[resist] (0,3) rectangle (1.5,3.6);
\fill[resist] (3,3) rectangle (12,3.6);
\fill[resist] (13.5,3) rectangle (15,3.6);
\fill[resist] (16.5,3) rectangle (20,3.6);
	\end{tikzpicture}
	\begin{tikzpicture}[node distance = 3cm, auto, thick,scale=\CrossAndTopSection, every node/.style={transform shape}]
		\fill[orange] (0,0) rectangle (20,12);

% n+
\fill[gray] (1.5,2) rectangle (3,6.5);
\fill[gray] (12,2) rectangle (16.5,6.5);

	\end{tikzpicture}
	\drawStepArrow{}
	\begin{tikzpicture}[node distance = 3cm, auto, thick,scale=\CrossAndTopSection, every node/.style={transform shape}]
		\input{tikz_process_steps/well.a.tex}
% gate oxide
\fill[LightGray] (4.8,2) rectangle (6.7,2.3);
\fill[LightGray] (13.3,2) rectangle (15.2,2.3);
% gate poly
\fill[BrickRed] (4.8,2.3) rectangle (6.7,2.6);
\fill[BrickRed] (13.3,2.3) rectangle (15.2,2.6);
% oxide
\fill[gray] (0,2) rectangle (1.5,3);
\fill[gray] (3,2) rectangle (4.8,3);
\fill[gray] (4.8,2.6) rectangle (6.7,3);
\fill[gray] (6.7,2) rectangle (12,3);
\fill[gray] (16.5,2) rectangle (20,3);

% resist
\fill[orange] (0,3) rectangle (1.5,3.6);
\fill[orange] (3,3) rectangle (12,3.6);
\fill[orange] (16.5,3) rectangle (20,3.6);;
	\end{tikzpicture}
	\begin{tikzpicture}[node distance = 3cm, auto, thick,scale=\CrossAndTopSection, every node/.style={transform shape}]
		\fill[resist] (0,0) rectangle (20,12);

% n+
\fill[nwell] (1.5,2) rectangle (3,6.5);
\fill[substrate] (12,2) rectangle (13.5,6.5);
\fill[substrate] (15,2) rectangle (16.5,6.5);
	\end{tikzpicture}
	\caption{N+ region opened}
\end{figure}

\subsection{Cleaning}
\begin{figure}[H]
	\centering
	\begin{tikzpicture}[node distance = 3cm, auto, thick,scale=\CrossSectionOnly, every node/.style={transform shape}]
		\input{tikz_process_steps/nimplant.4.a.tex}
	\end{tikzpicture}
	\drawStepArrow{}
	\begin{tikzpicture}[node distance = 3cm, auto, thick,scale=\CrossSectionOnly, every node/.style={transform shape}]
		\input{tikz_process_steps/well.a.tex}
% oxide
\fill[gray] (0,2) rectangle (3.5,3);
\fill[gray] (5,2) rectangle (6.5,3);
\fill[gray] (8,2) rectangle (17,3);
\fill[gray] (18.5,2) rectangle (20,3);
	\end{tikzpicture}
	\caption{Resist removal}
\end{figure}

\subsection{Implantation/Doping}\label{nimplant_implant_step}

We now need to bring in the carriers in order to build the n-junctions.

\begin{figure}[H]
	\centering
	\begin{tikzpicture}[node distance = 3cm, auto, thick,scale=\CrossSectionOnly, every node/.style={transform shape}]
		\input{tikz_process_steps/well.a.tex}
% oxide
\fill[isolationoxide] (0,2) rectangle (1.5,3);
\fill[isolationoxide] (3,2) rectangle (12,3);
\fill[isolationoxide] (13.5,2) rectangle (15,3);
\fill[isolationoxide] (16.5,2) rectangle (20,3);

\forloop{ct}{0}{\value{ct} < 21}
{
	\draw [->] (\value{ct},4) -- (\value{ct},3);
	\node at (\value{ct},4.2) {P$^{31}$};
}
	\end{tikzpicture}
	\drawStepArrow{}
	\begin{tikzpicture}[node distance = 3cm, auto, thick,scale=\CrossSectionOnly, every node/.style={transform shape}]
		\input{tikz_process_steps/nwell.a.tex}
% p-well
\fill[pwell] (11.5,0.75) rectangle (19,2);
\node at (14.25,1) {P-Well};
% oxide
\fill[isolationoxide] (0,2) rectangle (1.5,3);
\fill[isolationoxide] (3,2) rectangle (12,3);
\fill[isolationoxide] (13.5,2) rectangle (15,3);
\fill[isolationoxide] (16.5,2) rectangle (20,3);

\fill[nimplant] (1.5,1) rectangle (3,2);
\node at (2,1.5) {n+};
\fill[nimplant] (15,1) rectangle (16.5,2);
\node at (16,1.5) {n+};
\fill[nimplant] (12,1) rectangle (13.5,2);
\node at (13,1.5) {n+};
	\end{tikzpicture}
	\caption{N+ injection process}
\end{figure}

\textbf{Possible approaches}:
\begin{itemize}
	\item \textbf{"CF-3000 Implanter (IMP-3000)" from HKUST} \\
	At HKUST we have an implanter which gives us better control over the initial surface concentration. \\
	These steps are needed to arrive with the desired geometry:
	\begin{enumerate}
		\item Preparing by default cleaning
		\item The N-well is implanted with a Phosphorus ($P^{31}$) dose of $2.5\times10^{12}cm^{-2}$ at an energy of 100 keV.
		\item The N-well is annealed for 30 minutes at 1050\degreesC in $N_2$ environment (DIF-A1)\\
		After that the P-well will be around 2\um deep and the N-Well around 1\um deep
	\end{enumerate}
	\item \textbf{Constant source diffusion} \\
	We can add a layer of Phosphorus solution and diffusing in order to have an initial concentration in order to reach the desired concentration later by main diffusion.
\end{itemize}

\subsection{Oxide removal}
\begin{figure}[H]
	\centering
	\begin{tikzpicture}[node distance = 3cm, auto, thick,scale=\CrossSectionOnly, every node/.style={transform shape}]
		\input{tikz_process_steps/well.a.tex}
% gate oxide
\fill[LightGray] (4.8,2) rectangle (6.7,2.3);
\fill[LightGray] (13.3,2) rectangle (15.2,2.3);
% gate poly
\fill[BrickRed] (4.8,2.3) rectangle (6.7,2.6);
\fill[BrickRed] (13.3,2.3) rectangle (15.2,2.6);
% oxide
\fill[gray] (0,2) rectangle (1.5,3);
\fill[gray] (3,2) rectangle (4.8,3);
\fill[gray] (4.8,2.6) rectangle (6.7,3);
\fill[gray] (6.7,2) rectangle (12,3);
\fill[gray] (16.5,2) rectangle (20,3);

\fill[ProcessBlue] (1.5,1) rectangle (3,2);
\node at (2,1.5) {n+};
\fill[ProcessBlue] (15,1) rectangle (16.5,2);
\node at (16,1.5) {n+};
\fill[ProcessBlue] (12,1) rectangle (13.5,2);
\node at (13,1.5) {n+};
	\end{tikzpicture}
	\drawStepArrow{}
	\begin{tikzpicture}[node distance = 3cm, auto, thick,scale=\CrossSectionOnly, every node/.style={transform shape}]
		\input{tikz_process_steps/well.a.tex}
\fill[nimplant] (1.5,1) rectangle (3,2);
\node at (2,1.5) {n+};
\fill[nimplant] (15,1) rectangle (16.5,2);
\node at (16,1.5) {n+};
\fill[nimplant] (12,1) rectangle (13.5,2);
\node at (13,1.5) {n+};
	\end{tikzpicture}
	\caption{Oxide removal}
\end{figure}

\newpage
\section{p+ Implant}\label{pimplant}
For the bulk of the NMOS transistors and for the source and drain of the PMOS transistors highly doped  p+ areas are required.
In this step we're going to build these.

\begin{figure}[H]
	\centering
	\begin{tikzpicture}[node distance = 3cm, auto, thick,scale=\CrossAndTopSectionBig, every node/.style={transform shape}]
		% substrate
\fill[YellowOrange] (0,0) rectangle (20,2);
\node at (2,0.5) {Si (p-type)};
% n-well
\fill[Goldenrod] (1.25,0.75) rectangle (8.25,2);
\node at (5.75,1) {N-Well};
% body
\fill[ProcessBlue] (1.5,1) rectangle (3,2);
\node at (2,1.5) {n+};
% source
\fill[RedOrange] (3.5,1) rectangle (5,2);
\node at (4,1.5) {p+};
% drain
\fill[RedOrange] (6.5,1) rectangle (8,2);
\node at (7,1.5) {p+};
%% gate:
% gate oxide
\fill[LightGray] (4.8,2) rectangle (6.7,2.1);
% gate poly
\fill[BrickRed] (4.8,2.1) rectangle (6.7,2.2);

%field oxides:
\fill[DarkGray] (0,2) rectangle (1,4);
\fill[DarkGray] (8.5,2) rectangle (11.5,4);
\fill[DarkGray] (19,2) rectangle (20,4);

\fill[RedOrange] (0,1.5) rectangle (1,2);
\fill[RedOrange] (8.5,1.5) rectangle (11.5,2);
\fill[RedOrange] (19,1.5) rectangle (20,2);

\node at (0.5,1.75) {p+};
\node at (9.5,1.75) {p+};
\node at (19.5,1.75) {p+};

%%% nmos:
% body
\fill[RedOrange] (17,1) rectangle (18.5,2);
\node at (18,1.5) {p+};
% source
\fill[ProcessBlue] (15,1) rectangle (16.5,2);
\node at (16,1.5) {n+};
% drain
\fill[ProcessBlue] (12,1) rectangle (13.5,2);
\node at (13,1.5) {n+};

%% gate:
% gate oxide
\fill[LightGray] (13.3,2) rectangle (15.2,2.1);
% gate poly
\fill[BrickRed] (13.3,2.1) rectangle (15.2,2.2);
\fill[pimplant] (3.5,1) rectangle (5,2);
\node at (4,1.5) {p+};
\fill[pimplant] (6.5,1) rectangle (8,2);
\node at (7,1.5) {p+};
\fill[pimplant] (17,1) rectangle (18.5,2);
\node at (18,1.5) {p+};
	\end{tikzpicture}
	\begin{tikzpicture}[node distance = 3cm, auto, thick,scale=\CrossAndTopSectionBig, every node/.style={transform shape}]
		\fill[YellowOrange] (0,0) rectangle (20,10);

% n-well
\fill[Goldenrod] (1,1.25) rectangle (8.5,7.5);

% p+
\fill[RedOrange] (3.5,2) rectangle (5,6.5);
\fill[RedOrange] (6.5,2) rectangle (8,6.5);
\fill[RedOrange] (17,2) rectangle (18.5,6.5);

% n+
\fill[ProcessBlue] (1.5,2) rectangle (3,6.5);
\fill[ProcessBlue] (12,2) rectangle (13.5,6.5);
\fill[ProcessBlue] (15,2) rectangle (16.5,6.5);

% trench area
\fill[DarkGray] (0,0) rectangle (1,12);
\fill[DarkGray] (8.5,0) rectangle (11.5,12);
\fill[DarkGray] (19,0) rectangle (20,12);
\fill[DarkGray] (0,0) rectangle (20,1.25);
\fill[DarkGray] (0,7.5) rectangle (20,12);

% poly
\fill[BrickRed] (4.8,1.75) rectangle (6.7,9);
\fill[BrickRed] (13.3,1.75) rectangle (15.2,9);
\fill[BrickRed] (4.8,8) rectangle (15.2,9);
	\end{tikzpicture}
	\caption{P+ implant geometry target}
\end{figure}

The tricky thing here is to have a reasonable implant depth but not too deep because the deeper the junction, the higher the junction capacity which in turn limits the switching performance of the CMOS circuitry.

\begin{figure}[H]
	\centering
	\begin{tikzpicture}[node distance =1cm, auto, thick,scale=\VLSILayout, every node/.style={transform shape}]
		\input{tikz_process_steps/gate.layout.tex}

% n+
\fill[nimplant,opacity=\OpacityLayout] (1.5,2) rectangle (3,6.5);
\fill[nimplant,opacity=\OpacityLayout] (12,2) rectangle (13.5,6.5);
\fill[nimplant,opacity=\OpacityLayout] (15,2) rectangle (16.5,6.5);


% p+
\fill[pimplant,opacity=\OpacityLayout] (3.5,2) rectangle (5,6.5);
\fill[pimplant,opacity=\OpacityLayout] (6.5,2) rectangle (8,6.5);
\fill[pimplant,opacity=\OpacityLayout] (17,2) rectangle (18.5,6.5);
	\end{tikzpicture}
	\caption{P+ layout}
	\label{pimplant_layout}
\end{figure}

An example layout of p-implants can be seen in \autoref{pimplant_layout}, the mask is being extracted from the layer "p\_plus\_select"

Also important to notice is that this example layout is just for demonstration purposes only, please have a look at the standard cell documentation for the actual layouts. 

\newpage

\subsection{Mask dioxide layer}
\begin{figure}[H]
	\centering
	\begin{tikzpicture}[node distance = 3cm, auto, thick,scale=\CrossSectionOnly, every node/.style={transform shape}]
		\input{tikz_process_steps/pimplant.oxide_growth.a.tex}
	\end{tikzpicture}
	\drawStepArrow{}
	\begin{tikzpicture}[node distance = 3cm, auto, thick,scale=\CrossSectionOnly, every node/.style={transform shape}]
		% oxide
\fill[isolationoxide] (0,2.0) rectangle (20,4.0);

\fill[isolationoxide] (0,4.0) rectangle (0.5,4.75);
\fill[isolationoxide] (9.00,4.0) rectangle (11.0,4.75);
\fill[isolationoxide] (19.5,4.0) rectangle (20,4.75);

\filldraw[line width=0, isolationoxide] (0.5,4.75) -- (0.5,4.0) -- (1.25,4.0);
\filldraw[line width=0, isolationoxide] (8.5,4.0) -- (9.00,4.0) -- (9.00,4.75);

\filldraw[line width=0, isolationoxide] (11.0,4.75) -- (11.0,4.0) -- (11.75,4.0);
\filldraw[line width=0, isolationoxide] (18.75,4.0) -- (19.5,4.0) -- (19.5,4.75);

% substrate
\fill[YellowOrange] (0,0) rectangle (20,2);
\node at (2,0.5) {Si (p-type)};
% n-well
\fill[Goldenrod] (1.25,0.75) rectangle (8.25,2);
\node at (5.75,1) {N-Well};
% body
\fill[ProcessBlue] (1.5,1) rectangle (3,2);
\node at (2,1.5) {n+};
% source
\fill[RedOrange] (3.5,1) rectangle (5,2);
\node at (4,1.5) {p+};
% drain
\fill[RedOrange] (6.5,1) rectangle (8,2);
\node at (7,1.5) {p+};
%% gate:
% gate oxide
\fill[LightGray] (4.8,2) rectangle (6.7,2.1);
% gate poly
\fill[BrickRed] (4.8,2.1) rectangle (6.7,2.2);

%field oxides:
\fill[DarkGray] (0,2) rectangle (1,4);
\fill[DarkGray] (8.5,2) rectangle (11.5,4);
\fill[DarkGray] (19,2) rectangle (20,4);

\fill[RedOrange] (0,1.5) rectangle (1,2);
\fill[RedOrange] (8.5,1.5) rectangle (11.5,2);
\fill[RedOrange] (19,1.5) rectangle (20,2);

\node at (0.5,1.75) {p+};
\node at (9.5,1.75) {p+};
\node at (19.5,1.75) {p+};

%%% nmos:
% body
\fill[RedOrange] (17,1) rectangle (18.5,2);
\node at (18,1.5) {p+};
% source
\fill[ProcessBlue] (15,1) rectangle (16.5,2);
\node at (16,1.5) {n+};
% drain
\fill[ProcessBlue] (12,1) rectangle (13.5,2);
\node at (13,1.5) {n+};

%% gate:
% gate oxide
\fill[LightGray] (13.3,2) rectangle (15.2,2.1);
% gate poly
\fill[BrickRed] (13.3,2.1) rectangle (15.2,2.2);
	\end{tikzpicture}
	\caption{Oxide layer}
\end{figure}

\subsection{Pattering}
\begin{figure}[H]
	\centering
	\begin{tikzpicture}[node distance = 3cm, auto, thick,scale=\CrossAndTopSection, every node/.style={transform shape}]
		% substrate
\fill[YellowOrange] (0,0) rectangle (20,2);
\node at (2,0.5) {Si (p-type)};
% n-well
\fill[Goldenrod] (1.25,0.75) rectangle (8.25,2);
\node at (5.75,1) {N-Well};
% body
\fill[ProcessBlue] (1.5,1) rectangle (3,2);
\node at (2,1.5) {n+};
% source
\fill[RedOrange] (3.5,1) rectangle (5,2);
\node at (4,1.5) {p+};
% drain
\fill[RedOrange] (6.5,1) rectangle (8,2);
\node at (7,1.5) {p+};
%% gate:
% gate oxide
\fill[LightGray] (4.8,2) rectangle (6.7,2.1);
% gate poly
\fill[BrickRed] (4.8,2.1) rectangle (6.7,2.2);

%field oxides:
\fill[DarkGray] (0,2) rectangle (1,4);
\fill[DarkGray] (8.5,2) rectangle (11.5,4);
\fill[DarkGray] (19,2) rectangle (20,4);

\fill[RedOrange] (0,1.5) rectangle (1,2);
\fill[RedOrange] (8.5,1.5) rectangle (11.5,2);
\fill[RedOrange] (19,1.5) rectangle (20,2);

\node at (0.5,1.75) {p+};
\node at (9.5,1.75) {p+};
\node at (19.5,1.75) {p+};

%%% nmos:
% body
\fill[RedOrange] (17,1) rectangle (18.5,2);
\node at (18,1.5) {p+};
% source
\fill[ProcessBlue] (15,1) rectangle (16.5,2);
\node at (16,1.5) {n+};
% drain
\fill[ProcessBlue] (12,1) rectangle (13.5,2);
\node at (13,1.5) {n+};

%% gate:
% gate oxide
\fill[LightGray] (13.3,2) rectangle (15.2,2.1);
% gate poly
\fill[BrickRed] (13.3,2.1) rectangle (15.2,2.2);

	\end{tikzpicture}
	\begin{tikzpicture}[node distance = 3cm, auto, thick,scale=\CrossAndTopSection, every node/.style={transform shape}]
		\input{tikz_process_steps/pimplant.patterning.at.tex}
	\end{tikzpicture}
	\drawStepArrow{}
	\begin{tikzpicture}[node distance = 3cm, auto, thick,scale=\CrossAndTopSection, every node/.style={transform shape}]
		% resist
\fill[resist] (0,2.0) rectangle (3.75,5.0);
\fill[resist] (8.75,2.0) rectangle (16.75,5.0);
\fill[resist] (19.25,2.0) rectangle (20,5.0);

\input{tikz_process_steps/nimplant.a.tex}


	\end{tikzpicture}
	\begin{tikzpicture}[node distance = 3cm, auto, thick,scale=\CrossAndTopSection, every node/.style={transform shape}]
		\fill[resist] (0,0) rectangle (20,12);

% p+
\fill[isolationoxide] (3.5,2) rectangle (5,6.5);
\fill[isolationoxide] (6.5,2) rectangle (8,6.5);
\fill[isolationoxide] (17,2) rectangle (18.25,6.5);
	\end{tikzpicture}
	\caption{P+ region resist mask}
\end{figure}

\subsection{Etching}
\begin{figure}[H]
	\centering
	\begin{tikzpicture}[node distance = 3cm, auto, thick,scale=\CrossAndTopSection, every node/.style={transform shape}]
		% oxide
\fill[isolationoxide] (0,2) rectangle (20,3.5);
% resist
\fill[resist] (0,3.5) rectangle (3.5,4.1);
\fill[resist] (8,3.5) rectangle (17,4.1);
\fill[resist] (18.5,3.5) rectangle (20,4.1);

% substrate
\fill[YellowOrange] (0,0) rectangle (20,2);
\node at (2,0.5) {Si (p-type)};
% n-well
\fill[Goldenrod] (1.25,0.75) rectangle (8.25,2);
\node at (5.75,1) {N-Well};
% body
\fill[ProcessBlue] (1.5,1) rectangle (3,2);
\node at (2,1.5) {n+};
% source
\fill[RedOrange] (3.5,1) rectangle (5,2);
\node at (4,1.5) {p+};
% drain
\fill[RedOrange] (6.5,1) rectangle (8,2);
\node at (7,1.5) {p+};
%% gate:
% gate oxide
\fill[LightGray] (4.8,2) rectangle (6.7,2.1);
% gate poly
\fill[BrickRed] (4.8,2.1) rectangle (6.7,2.2);

%field oxides:
\fill[DarkGray] (0,2) rectangle (1,4);
\fill[DarkGray] (8.5,2) rectangle (11.5,4);
\fill[DarkGray] (19,2) rectangle (20,4);

\fill[RedOrange] (0,1.5) rectangle (1,2);
\fill[RedOrange] (8.5,1.5) rectangle (11.5,2);
\fill[RedOrange] (19,1.5) rectangle (20,2);

\node at (0.5,1.75) {p+};
\node at (9.5,1.75) {p+};
\node at (19.5,1.75) {p+};

%%% nmos:
% body
\fill[RedOrange] (17,1) rectangle (18.5,2);
\node at (18,1.5) {p+};
% source
\fill[ProcessBlue] (15,1) rectangle (16.5,2);
\node at (16,1.5) {n+};
% drain
\fill[ProcessBlue] (12,1) rectangle (13.5,2);
\node at (13,1.5) {n+};

%% gate:
% gate oxide
\fill[LightGray] (13.3,2) rectangle (15.2,2.1);
% gate poly
\fill[BrickRed] (13.3,2.1) rectangle (15.2,2.2);
	\end{tikzpicture}
	\begin{tikzpicture}[node distance = 3cm, auto, thick,scale=\CrossAndTopSection, every node/.style={transform shape}]
		\fill[orange] (0,0) rectangle (20,12);

% n+
\fill[gray] (3.5,2) rectangle (8,6.5);
\fill[gray] (17,2) rectangle (18.5,6.5);
	\end{tikzpicture}
	\drawStepArrow{}
	\begin{tikzpicture}[node distance = 3cm, auto, thick,scale=\CrossAndTopSection, every node/.style={transform shape}]
		% oxide
\fill[isolationoxide] (0,2) rectangle (3.0,3.5);
\fill[isolationoxide] (8.5,2) rectangle (17,3.5);
\fill[isolationoxide] (18.75,2) rectangle (20,3.5);

% resist
\fill[resist] (0,3.5) rectangle (3.0,4.1);
\fill[resist] (8.5,3.5) rectangle (17,4.1);
\fill[resist] (18.75,3.5) rectangle (20,4.1);

% substrate
\fill[YellowOrange] (0,0) rectangle (20,2);
\node at (2,0.5) {Si (p-type)};
% n-well
\fill[Goldenrod] (1.25,0.75) rectangle (8.25,2);
\node at (5.75,1) {N-Well};
% body
\fill[ProcessBlue] (1.5,1) rectangle (3,2);
\node at (2,1.5) {n+};
% source
\fill[RedOrange] (3.5,1) rectangle (5,2);
\node at (4,1.5) {p+};
% drain
\fill[RedOrange] (6.5,1) rectangle (8,2);
\node at (7,1.5) {p+};
%% gate:
% gate oxide
\fill[LightGray] (4.8,2) rectangle (6.7,2.1);
% gate poly
\fill[BrickRed] (4.8,2.1) rectangle (6.7,2.2);

%field oxides:
\fill[DarkGray] (0,2) rectangle (1,4);
\fill[DarkGray] (8.5,2) rectangle (11.5,4);
\fill[DarkGray] (19,2) rectangle (20,4);

\fill[RedOrange] (0,1.5) rectangle (1,2);
\fill[RedOrange] (8.5,1.5) rectangle (11.5,2);
\fill[RedOrange] (19,1.5) rectangle (20,2);

\node at (0.5,1.75) {p+};
\node at (9.5,1.75) {p+};
\node at (19.5,1.75) {p+};

%%% nmos:
% body
\fill[RedOrange] (17,1) rectangle (18.5,2);
\node at (18,1.5) {p+};
% source
\fill[ProcessBlue] (15,1) rectangle (16.5,2);
\node at (16,1.5) {n+};
% drain
\fill[ProcessBlue] (12,1) rectangle (13.5,2);
\node at (13,1.5) {n+};

%% gate:
% gate oxide
\fill[LightGray] (13.3,2) rectangle (15.2,2.1);
% gate poly
\fill[BrickRed] (13.3,2.1) rectangle (15.2,2.2);

	\end{tikzpicture}
	\begin{tikzpicture}[node distance = 3cm, auto, thick,scale=\CrossAndTopSection, every node/.style={transform shape}]
		\fill[orange] (0,0) rectangle (20,12);

% p+
\fill[Goldenrod] (3.5,2) rectangle (5,6.5);
\fill[Goldenrod] (6.5,2) rectangle (8,6.5);
\fill[YellowOrange] (17,2) rectangle (18.5,6.5);
	\end{tikzpicture}
	\caption{P+ region opened}
\end{figure}

\subsection{Cleaning}
\begin{figure}[H]
	\centering
	\begin{tikzpicture}[node distance = 3cm, auto, thick,scale=\CrossSectionOnly, every node/.style={transform shape}]
		% substrate
\fill[YellowOrange] (0,0) rectangle (20,2);
\node at (2,0.5) {Si (p-type)};
% n-well
\fill[Goldenrod] (1.25,0.75) rectangle (8.25,2);
\node at (5.75,1) {N-Well};
% body
\fill[ProcessBlue] (1.5,1) rectangle (3,2);
\node at (2,1.5) {n+};
% source
\fill[RedOrange] (3.5,1) rectangle (5,2);
\node at (4,1.5) {p+};
% drain
\fill[RedOrange] (6.5,1) rectangle (8,2);
\node at (7,1.5) {p+};
%% gate:
% gate oxide
\fill[LightGray] (4.8,2) rectangle (6.7,2.1);
% gate poly
\fill[BrickRed] (4.8,2.1) rectangle (6.7,2.2);

%field oxides:
\fill[DarkGray] (0,2) rectangle (1,4);
\fill[DarkGray] (8.5,2) rectangle (11.5,4);
\fill[DarkGray] (19,2) rectangle (20,4);

\fill[RedOrange] (0,1.5) rectangle (1,2);
\fill[RedOrange] (8.5,1.5) rectangle (11.5,2);
\fill[RedOrange] (19,1.5) rectangle (20,2);

\node at (0.5,1.75) {p+};
\node at (9.5,1.75) {p+};
\node at (19.5,1.75) {p+};

%%% nmos:
% body
\fill[RedOrange] (17,1) rectangle (18.5,2);
\node at (18,1.5) {p+};
% source
\fill[ProcessBlue] (15,1) rectangle (16.5,2);
\node at (16,1.5) {n+};
% drain
\fill[ProcessBlue] (12,1) rectangle (13.5,2);
\node at (13,1.5) {n+};

%% gate:
% gate oxide
\fill[LightGray] (13.3,2) rectangle (15.2,2.1);
% gate poly
\fill[BrickRed] (13.3,2.1) rectangle (15.2,2.2);

% oxide
\fill[isolationoxide] (0,2) rectangle (3.5,3);
\fill[isolationoxide] (5,2) rectangle (6.5,3);
\fill[isolationoxide] (8,2) rectangle (17,3);
\fill[isolationoxide] (18.5,2) rectangle (20,3);;

% resist
\fill[resist] (0,3) rectangle (3.5,3.6);
\fill[resist] (5,3) rectangle (6.5,3.6);
\fill[resist] (8,3) rectangle (17,3.6);
\fill[resist] (18.5,3) rectangle (20,3.6);;
	\end{tikzpicture}
	\drawStepArrow{}
	\begin{tikzpicture}[node distance = 3cm, auto, thick,scale=\CrossSectionOnly, every node/.style={transform shape}]
		% substrate
\fill[YellowOrange] (0,0) rectangle (20,2);
\node at (2,0.5) {Si (p-type)};
% n-well
\fill[Goldenrod] (1.25,0.75) rectangle (8.25,2);
\node at (5.75,1) {N-Well};
% body
\fill[ProcessBlue] (1.5,1) rectangle (3,2);
\node at (2,1.5) {n+};
% source
\fill[RedOrange] (3.5,1) rectangle (5,2);
\node at (4,1.5) {p+};
% drain
\fill[RedOrange] (6.5,1) rectangle (8,2);
\node at (7,1.5) {p+};
%% gate:
% gate oxide
\fill[LightGray] (4.8,2) rectangle (6.7,2.1);
% gate poly
\fill[BrickRed] (4.8,2.1) rectangle (6.7,2.2);

%field oxides:
\fill[DarkGray] (0,2) rectangle (1,4);
\fill[DarkGray] (8.5,2) rectangle (11.5,4);
\fill[DarkGray] (19,2) rectangle (20,4);

\fill[RedOrange] (0,1.5) rectangle (1,2);
\fill[RedOrange] (8.5,1.5) rectangle (11.5,2);
\fill[RedOrange] (19,1.5) rectangle (20,2);

\node at (0.5,1.75) {p+};
\node at (9.5,1.75) {p+};
\node at (19.5,1.75) {p+};

%%% nmos:
% body
\fill[RedOrange] (17,1) rectangle (18.5,2);
\node at (18,1.5) {p+};
% source
\fill[ProcessBlue] (15,1) rectangle (16.5,2);
\node at (16,1.5) {n+};
% drain
\fill[ProcessBlue] (12,1) rectangle (13.5,2);
\node at (13,1.5) {n+};

%% gate:
% gate oxide
\fill[LightGray] (13.3,2) rectangle (15.2,2.1);
% gate poly
\fill[BrickRed] (13.3,2.1) rectangle (15.2,2.2);

% oxide
\fill[gray] (0,2) rectangle (3.5,3);
\fill[gray] (8,2) rectangle (13.3,3);
\fill[gray] (13.3,2.6) rectangle (15.2,3);
\fill[gray] (15.2,2) rectangle (17,3);
\fill[gray] (18.5,2) rectangle (20,3);
	\end{tikzpicture}
	\caption{Resist removal}
\end{figure}

\subsection{Implantation/Doping}\label{pimplant_implant_step}
\begin{figure}[H]
	\centering
	\begin{tikzpicture}[node distance = 3cm, auto, thick,scale=\CrossSectionOnly, every node/.style={transform shape}]
		% oxide
\fill[isolationoxide] (0,2) rectangle (3.5,3.5);
\fill[isolationoxide] (8,2) rectangle (17,3.5);
\fill[isolationoxide] (18.5,2) rectangle (20,3.5);

\forloop{ct}{0}{\value{ct} < 21}
{
	\draw [->] (\value{ct},5) -- (\value{ct},4);
	\node at (\value{ct},5.2) {B$^{11}$};
}

% substrate
\fill[YellowOrange] (0,0) rectangle (20,2);
\node at (2,0.5) {Si (p-type)};
% n-well
\fill[Goldenrod] (1.25,0.75) rectangle (8.25,2);
\node at (5.75,1) {N-Well};
% body
\fill[ProcessBlue] (1.5,1) rectangle (3,2);
\node at (2,1.5) {n+};
% source
\fill[RedOrange] (3.5,1) rectangle (5,2);
\node at (4,1.5) {p+};
% drain
\fill[RedOrange] (6.5,1) rectangle (8,2);
\node at (7,1.5) {p+};
%% gate:
% gate oxide
\fill[LightGray] (4.8,2) rectangle (6.7,2.1);
% gate poly
\fill[BrickRed] (4.8,2.1) rectangle (6.7,2.2);

%field oxides:
\fill[DarkGray] (0,2) rectangle (1,4);
\fill[DarkGray] (8.5,2) rectangle (11.5,4);
\fill[DarkGray] (19,2) rectangle (20,4);

\fill[RedOrange] (0,1.5) rectangle (1,2);
\fill[RedOrange] (8.5,1.5) rectangle (11.5,2);
\fill[RedOrange] (19,1.5) rectangle (20,2);

\node at (0.5,1.75) {p+};
\node at (9.5,1.75) {p+};
\node at (19.5,1.75) {p+};

%%% nmos:
% body
\fill[RedOrange] (17,1) rectangle (18.5,2);
\node at (18,1.5) {p+};
% source
\fill[ProcessBlue] (15,1) rectangle (16.5,2);
\node at (16,1.5) {n+};
% drain
\fill[ProcessBlue] (12,1) rectangle (13.5,2);
\node at (13,1.5) {n+};

%% gate:
% gate oxide
\fill[LightGray] (13.3,2) rectangle (15.2,2.1);
% gate poly
\fill[BrickRed] (13.3,2.1) rectangle (15.2,2.2);

	\end{tikzpicture}
	\drawStepArrow{}
	\begin{tikzpicture}[node distance = 3cm, auto, thick,scale=\CrossSectionOnly, every node/.style={transform shape}]
		% substrate
\fill[YellowOrange] (0,0) rectangle (20,2);
\node at (2,0.5) {Si (p-type)};
% n-well
\fill[Goldenrod] (1.25,0.75) rectangle (8.25,2);
\node at (5.75,1) {N-Well};
% body
\fill[ProcessBlue] (1.5,1) rectangle (3,2);
\node at (2,1.5) {n+};
% source
\fill[RedOrange] (3.5,1) rectangle (5,2);
\node at (4,1.5) {p+};
% drain
\fill[RedOrange] (6.5,1) rectangle (8,2);
\node at (7,1.5) {p+};
%% gate:
% gate oxide
\fill[LightGray] (4.8,2) rectangle (6.7,2.1);
% gate poly
\fill[BrickRed] (4.8,2.1) rectangle (6.7,2.2);

%field oxides:
\fill[DarkGray] (0,2) rectangle (1,4);
\fill[DarkGray] (8.5,2) rectangle (11.5,4);
\fill[DarkGray] (19,2) rectangle (20,4);

\fill[RedOrange] (0,1.5) rectangle (1,2);
\fill[RedOrange] (8.5,1.5) rectangle (11.5,2);
\fill[RedOrange] (19,1.5) rectangle (20,2);

\node at (0.5,1.75) {p+};
\node at (9.5,1.75) {p+};
\node at (19.5,1.75) {p+};

%%% nmos:
% body
\fill[RedOrange] (17,1) rectangle (18.5,2);
\node at (18,1.5) {p+};
% source
\fill[ProcessBlue] (15,1) rectangle (16.5,2);
\node at (16,1.5) {n+};
% drain
\fill[ProcessBlue] (12,1) rectangle (13.5,2);
\node at (13,1.5) {n+};

%% gate:
% gate oxide
\fill[LightGray] (13.3,2) rectangle (15.2,2.1);
% gate poly
\fill[BrickRed] (13.3,2.1) rectangle (15.2,2.2);
% oxide
\fill[isolationoxide] (0,2) rectangle (3.5,3);
\fill[isolationoxide] (5,2) rectangle (6.5,3);
\fill[isolationoxide] (8,2) rectangle (17,3);
\fill[isolationoxide] (18.5,2) rectangle (20,3);

\fill[pimplant] (3.5,1) rectangle (5,2);
\node at (4,1.5) {p+};
\fill[pimplant] (6.5,1) rectangle (8,2);
\node at (7,1.5) {p+};
\fill[pimplant] (17,1) rectangle (18.5,2);
\node at (18,1.5) {p+};
	\end{tikzpicture}
	\caption{P+ injection process}
\end{figure}

\subsection{Oxide removal}
\begin{figure}[H]
	\centering
	\begin{tikzpicture}[node distance = 3cm, auto, thick,scale=\CrossSectionOnly, every node/.style={transform shape}]
		% oxide
\fill[isolationoxide] (0,2) rectangle (3.0,3.5);
\fill[isolationoxide] (8.5,2) rectangle (17,3.5);
\fill[isolationoxide] (18.75,2) rectangle (20,3.5);

% substrate
\fill[YellowOrange] (0,0) rectangle (20,2);
\node at (2,0.5) {Si (p-type)};
% n-well
\fill[Goldenrod] (1.25,0.75) rectangle (8.25,2);
\node at (5.75,1) {N-Well};
% body
\fill[ProcessBlue] (1.5,1) rectangle (3,2);
\node at (2,1.5) {n+};
% source
\fill[RedOrange] (3.5,1) rectangle (5,2);
\node at (4,1.5) {p+};
% drain
\fill[RedOrange] (6.5,1) rectangle (8,2);
\node at (7,1.5) {p+};
%% gate:
% gate oxide
\fill[LightGray] (4.8,2) rectangle (6.7,2.1);
% gate poly
\fill[BrickRed] (4.8,2.1) rectangle (6.7,2.2);

%field oxides:
\fill[DarkGray] (0,2) rectangle (1,4);
\fill[DarkGray] (8.5,2) rectangle (11.5,4);
\fill[DarkGray] (19,2) rectangle (20,4);

\fill[RedOrange] (0,1.5) rectangle (1,2);
\fill[RedOrange] (8.5,1.5) rectangle (11.5,2);
\fill[RedOrange] (19,1.5) rectangle (20,2);

\node at (0.5,1.75) {p+};
\node at (9.5,1.75) {p+};
\node at (19.5,1.75) {p+};

%%% nmos:
% body
\fill[RedOrange] (17,1) rectangle (18.5,2);
\node at (18,1.5) {p+};
% source
\fill[ProcessBlue] (15,1) rectangle (16.5,2);
\node at (16,1.5) {n+};
% drain
\fill[ProcessBlue] (12,1) rectangle (13.5,2);
\node at (13,1.5) {n+};

%% gate:
% gate oxide
\fill[LightGray] (13.3,2) rectangle (15.2,2.1);
% gate poly
\fill[BrickRed] (13.3,2.1) rectangle (15.2,2.2);

\fill[pimplant] (3.0,1.5) rectangle (5,2);
\fill[pimplant] (6.5,1.5) rectangle (8.5,2);
\fill[pimplant] (17,1.5) rectangle (18.75,2);
	\end{tikzpicture}
	\drawStepArrow{}
	\begin{tikzpicture}[node distance = 3cm, auto, thick,scale=\CrossSectionOnly, every node/.style={transform shape}]
		% substrate
\fill[YellowOrange] (0,0) rectangle (20,2);
\node at (2,0.5) {Si (p-type)};
% n-well
\fill[Goldenrod] (1.25,0.75) rectangle (8.25,2);
\node at (5.75,1) {N-Well};
% body
\fill[ProcessBlue] (1.5,1) rectangle (3,2);
\node at (2,1.5) {n+};
% source
\fill[RedOrange] (3.5,1) rectangle (5,2);
\node at (4,1.5) {p+};
% drain
\fill[RedOrange] (6.5,1) rectangle (8,2);
\node at (7,1.5) {p+};
%% gate:
% gate oxide
\fill[LightGray] (4.8,2) rectangle (6.7,2.1);
% gate poly
\fill[BrickRed] (4.8,2.1) rectangle (6.7,2.2);

%field oxides:
\fill[DarkGray] (0,2) rectangle (1,4);
\fill[DarkGray] (8.5,2) rectangle (11.5,4);
\fill[DarkGray] (19,2) rectangle (20,4);

\fill[RedOrange] (0,1.5) rectangle (1,2);
\fill[RedOrange] (8.5,1.5) rectangle (11.5,2);
\fill[RedOrange] (19,1.5) rectangle (20,2);

\node at (0.5,1.75) {p+};
\node at (9.5,1.75) {p+};
\node at (19.5,1.75) {p+};

%%% nmos:
% body
\fill[RedOrange] (17,1) rectangle (18.5,2);
\node at (18,1.5) {p+};
% source
\fill[ProcessBlue] (15,1) rectangle (16.5,2);
\node at (16,1.5) {n+};
% drain
\fill[ProcessBlue] (12,1) rectangle (13.5,2);
\node at (13,1.5) {n+};

%% gate:
% gate oxide
\fill[LightGray] (13.3,2) rectangle (15.2,2.1);
% gate poly
\fill[BrickRed] (13.3,2.1) rectangle (15.2,2.2);

\fill[RedOrange] (3.5,1) rectangle (5,2);
\node at (4,1.5) {p+};
\fill[RedOrange] (6.5,1) rectangle (8,2);
\node at (7,1.5) {p+};
\fill[RedOrange] (17,1) rectangle (18.5,2);
\node at (18,1.5) {p+};
	\end{tikzpicture}
	\caption{Oxide removal}
\end{figure}

\newpage
\section{Gate}\label{gate}
Now we have to build the initial gate structure which contains of the 40nm thick dielectric (in our case just silicon dioxide) and the polysilicon electrode.

\begin{figure}[H]
	\centering
	\begin{tikzpicture}[node distance = 3cm, auto, thick,scale=\CrossAndTopSectionBig, every node/.style={transform shape}]
		\input{tikz_process_steps/nwell.a.tex}
% p-well
\fill[pwell] (11.5,0.75) rectangle (19,2);
\node at (14.25,1) {P-Well};
\fill[gateoxide] (4.8,2) rectangle (6.7,2.3);
\fill[gateoxide] (13.3,2) rectangle (15.2,2.3);
\fill[gatemetal] (4.8,2.3) rectangle (6.7,2.6);
\fill[gatemetal] (13.3,2.3) rectangle (15.2,2.6);
		\node at (3,3) {Gate oxide};
		\draw[->] (3,2.8) -- (5,2.2);
		\node at (3,4) {Polysilicon};
		\draw[->] (3,3.8) -- (5,3.2);
	\end{tikzpicture}
	\begin{tikzpicture}[node distance = 3cm, auto, thick,scale=\CrossAndTopSectionBig, every node/.style={transform shape}]
		\input{tikz_process_steps/nwell.b.tex}
\fill[pwell] (11.5,1.25) rectangle (19,7.5);

% gate metal
\fill[gatemetal] (5,0) rectangle (6.5,9);
\fill[gatemetal] (13.5,0) rectangle (15,9);
\fill[gatemetal] (5,8) rectangle (15,10);
	\end{tikzpicture}
	\caption{Poly silicon gate contacts with gate oxide}
\end{figure}

The line spacing of the polysilicon electrode shape has to be at least 0.5\um because of the resolution of the stepper and also because of the etching process which has 0.5\um as the minimum line spacing.

\begin{figure}[H]
	\centering
	\begin{tikzpicture}[node distance =1cm, auto, thick,scale=\VLSILayout, every node/.style={transform shape}]
		\input{tikz_process_steps/nimplant.layout.tex}

% p+
\fill[pimplant,opacity=\OpacityLayout] (3.5,2) rectangle (5,6.5);
\fill[pimplant,opacity=\OpacityLayout] (6.5,2) rectangle (8,6.5);
\fill[pimplant,opacity=\OpacityLayout] (17,2) rectangle (18.5,6.5);
% gate metal
\fill[gatemetal,opacity=\OpacityLayout] (4.8,1.75) rectangle (6.7,8);
\fill[gatemetal,opacity=\OpacityLayout] (13.3,1.75) rectangle (15.2,8);
\fill[gatemetal,opacity=\OpacityLayout] (4.8,8) rectangle (15.2,10);

	\end{tikzpicture}
	\caption{Gate layout}
	\label{gate_layout}
\end{figure}

In \autoref{gate_layout} we can see the layout honoring the 0.5\um spacing design rule for the gate structure shape and poly-layer interconnect between NMOS and PMOS.

\newpage

\subsection{Gate oxide deposition}

Now we have to deposit the dielectric isolator between the gate electrode and the channel.
As designed in the process design document, the layer will be 40nm thick.

\begin{figure}[H]
	\centering
	\begin{tikzpicture}[node distance = 3cm, auto, thick,scale=\CrossSectionOnly, every node/.style={transform shape}]
		\def\welldepthAfox{0.75}
\def\welldepthBfox{1.25}
\def\welldepthCfox{1.5}

\newcommand{\stopper}[1]{
\filldraw[line width=0, isolationoxide] (#1,2.0) -- (#1+0.75,2.0) -- (#1+0.75,2.75);
\fill[isolationoxide] (#1+0.75,2.0) rectangle (#1+1.25,2.75);
\filldraw[line width=0, isolationoxide] (#1+1.25,2.75) -- (#1+1.25,2.0) -- (#1+2.0,2.0);
}

\newcommand{\bjtstopper}[1]{
\filldraw[line width=0, isolationoxide] (#1,2.0) -- (#1+0.25,2.0) -- (#1+0.25,2.75);
\fill[isolationoxide] (#1+0.25,2.0) rectangle (#1+0.5,2.75);
\filldraw[line width=0, isolationoxide] (#1+0.5,2.75) -- (#1+0.5,2.0) -- (#1+0.75,2.0);
}

% oxide
\fill[isolationoxide] (0,1.25) rectangle (55.0,2.0);

\fill[isolationoxide] (0,2.0) rectangle (0.75,2.75);
\filldraw[line width=0, isolationoxide] (0.75,2.75) -- (0.75,2.0) -- (1.5,2.0);

%separates n+ from p+
\stopper{2.5}
\stopper{8.0}

%separates n+ from p+
\stopper{13.5}

\stopper{16.5}
\stopper{17.0}

\stopper{20.0}

\stopper{24.5}
\stopper{25.0}

\bjtstopper{27.75}
\bjtstopper{29.0}
\bjtstopper{30.5}
\bjtstopper{31.75}

\stopper{33.5}
\stopper{33.75}

\bjtstopper{36.5}
\bjtstopper{37.75}
\bjtstopper{39.0}
\bjtstopper{40.25}

\filldraw[line width=0, isolationoxide] (41.75,2.0) -- (42.5,2.0) -- (42.5,2.75);
\fill[isolationoxide] (42.5,2.0) rectangle (55.0,2.75);

\input{tikz_process_steps/sti.a.tex}

% substrate
\fill[substrate] (0,0) rectangle (55,1.25);
\node at (2,0.5) {Silicon substrate};

% normal wells
\shade[upper left = nwell, upper right = nwell, lower right = substrate, lower left = substrate,] (1.25,\welldepthAfox) rectangle (8.25,2.0);
\shade[upper left = pwell, upper right = pwell, lower right = substrate, lower left = substrate,] (9.75,\welldepthAfox) rectangle (16.75,2.0);
\shade[upper left = nwell, upper right = nwell, lower right = substrate, lower left = substrate,] (18.25,\welldepthAfox) rectangle (25.25,2.0);
\shade[upper left = nwell, upper right = nwell, lower right = substrate, lower left = substrate,] (26.75,\welldepthAfox) rectangle (33.75,2.0);
\shade[upper left = nwell, upper right = nwell, lower right = substrate, lower left = substrate,] (35.25,\welldepthAfox) rectangle (42.25,2.0);

% p base
\shade[upper left = pbase, upper right = pbase, lower right = nwell, lower left = nwell,] (18.5,\welldepthBfox) rectangle (25.0,2.0);

% npn - pbase
\shade[upper left = pbase, upper right = pbase, lower right = nwell, lower left = nwell,] (28.25,\welldepthBfox) rectangle (32.0,2.0);

% pnp emitter and collector ring
\shade[upper left = pbase, upper right = pbase, lower right = nwell, lower left = nwell,] (35.5,\welldepthBfox) rectangle (37.0,2.0);
\shade[upper left = pbase, upper right = pbase, lower right = nwell, lower left = nwell,] (38.0,\welldepthBfox) rectangle (39.5,2.0);
\shade[upper left = pbase, upper right = pbase, lower right = nwell, lower left = nwell,] (40.5,\welldepthBfox) rectangle (42.0,2.0);

% n base for SONOS
\shade[upper left = nbase, upper right = nbase, lower right = pbase, lower left = pbase,] (18.75,\welldepthCfox) rectangle (24.75,2.0);

% n base for NPN emitter
\shade[upper left = nbase, upper right = nbase, lower right = pbase, lower left = pbase,] (29.5,\welldepthCfox) rectangle (30.75,2.0);


	\end{tikzpicture} \\
	\includegraphics[scale=0.01]{down_arrow.png} \\
	\begin{tikzpicture}[node distance = 3cm, auto, thick,scale=\CrossSectionOnly, every node/.style={transform shape}]
		\def\welldepthAfox{0.75}
\def\welldepthBfox{1.25}
\def\welldepthCfox{1.5}

\newcommand{\stopper}[1]{
\filldraw[line width=0, isolationoxide] (#1,2.0) -- (#1+0.75,2.0) -- (#1+0.75,2.75);
\fill[isolationoxide] (#1+0.75,2.0) rectangle (#1+1.25,2.75);
\filldraw[line width=0, isolationoxide] (#1+1.25,2.75) -- (#1+1.25,2.0) -- (#1+2.0,2.0);
}

\newcommand{\bjtstopper}[1]{
\filldraw[line width=0, isolationoxide] (#1,2.0) -- (#1+0.25,2.0) -- (#1+0.25,2.75);
\fill[isolationoxide] (#1+0.25,2.0) rectangle (#1+0.5,2.75);
\filldraw[line width=0, isolationoxide] (#1+0.5,2.75) -- (#1+0.5,2.0) -- (#1+0.75,2.0);
}

% oxide
\fill[isolationoxide] (0,1.25) rectangle (55.0,2.0);

\fill[isolationoxide] (0,2.0) rectangle (0.75,2.75);
\filldraw[line width=0, isolationoxide] (0.75,2.75) -- (0.75,2.0) -- (1.5,2.0);

%separates n+ from p+
\stopper{2.5}
\stopper{8.0}

%separates n+ from p+
\stopper{13.5}

\stopper{16.5}
\stopper{17.0}

\stopper{20.0}

\stopper{24.5}
\stopper{25.0}

\bjtstopper{27.75}
\bjtstopper{29.0}
\bjtstopper{30.5}
\bjtstopper{31.75}

\stopper{33.5}
\stopper{33.75}

\bjtstopper{36.5}
\bjtstopper{37.75}
\bjtstopper{39.0}
\bjtstopper{40.25}

\filldraw[line width=0, isolationoxide] (41.75,2.0) -- (42.5,2.0) -- (42.5,2.75);
\fill[isolationoxide] (42.5,2.0) rectangle (55.0,2.75);

\input{tikz_process_steps/sti.a.tex}

% substrate
\fill[substrate] (0,0) rectangle (55,1.25);
\node at (2,0.5) {Silicon substrate};

% normal wells
\shade[upper left = nwell, upper right = nwell, lower right = substrate, lower left = substrate,] (1.25,\welldepthAfox) rectangle (8.25,2.0);
\shade[upper left = pwell, upper right = pwell, lower right = substrate, lower left = substrate,] (9.75,\welldepthAfox) rectangle (16.75,2.0);
\shade[upper left = nwell, upper right = nwell, lower right = substrate, lower left = substrate,] (18.25,\welldepthAfox) rectangle (25.25,2.0);
\shade[upper left = nwell, upper right = nwell, lower right = substrate, lower left = substrate,] (26.75,\welldepthAfox) rectangle (33.75,2.0);
\shade[upper left = nwell, upper right = nwell, lower right = substrate, lower left = substrate,] (35.25,\welldepthAfox) rectangle (42.25,2.0);

% p base
\shade[upper left = pbase, upper right = pbase, lower right = nwell, lower left = nwell,] (18.5,\welldepthBfox) rectangle (25.0,2.0);

% npn - pbase
\shade[upper left = pbase, upper right = pbase, lower right = nwell, lower left = nwell,] (28.25,\welldepthBfox) rectangle (32.0,2.0);

% pnp emitter and collector ring
\shade[upper left = pbase, upper right = pbase, lower right = nwell, lower left = nwell,] (35.5,\welldepthBfox) rectangle (37.0,2.0);
\shade[upper left = pbase, upper right = pbase, lower right = nwell, lower left = nwell,] (38.0,\welldepthBfox) rectangle (39.5,2.0);
\shade[upper left = pbase, upper right = pbase, lower right = nwell, lower left = nwell,] (40.5,\welldepthBfox) rectangle (42.0,2.0);

% n base for SONOS
\shade[upper left = nbase, upper right = nbase, lower right = pbase, lower left = pbase,] (18.75,\welldepthCfox) rectangle (24.75,2.0);

% n base for NPN emitter
\shade[upper left = nbase, upper right = nbase, lower right = pbase, lower left = pbase,] (29.5,\welldepthCfox) rectangle (30.75,2.0);


\fill[gateoxide] (0,2.75) rectangle (0.5,3.05);
\filldraw[line width=0, gateoxide] (0.5,2.75)--(1.25,2.0)--(1.25,2.3)--(0.5,3.05);
\fill[gateoxide] (1.25,2.0) rectangle (8.25,2.3);
\filldraw[line width=0, gateoxide] (9.00,2.75)--(8.25,2.0)--(8.25,2.3)--(9.00,3.05);
\fill[gateoxide] (9.00,2.75) rectangle (11.0,3.05);
\filldraw[line width=0, gateoxide] (11.0,2.75)--(11.75,2.0)--(11.75,2.3)--(11.0,3.05);
\fill[gateoxide] (11.75,2.0) rectangle (18.75,2.3);
\filldraw[line width=0, gateoxide] (19.5,2.75)--(18.75,2.0)--(18.75,2.3)--(19.5,3.05);
\fill[gateoxide] (19.5,2.75) rectangle (20.0,3.05);
	\end{tikzpicture}
	\caption{Thin oxide}
\end{figure}
The thickness of this layer decides over many critical key properties of the transistor, hence there should be little to no variation in the thickness of the gate oxide layer.
For that reason we put the wafer into the diffusion furnace and perform dry oxidation at 1050\degreesC for 33 minutes and 14 seconds.\footnote{\url{http://cleanroom.byu.edu/OxideTimeCalc}}

\subsection{Polysilicon deposition}

Now we need to add the polysilicon layer for forming the gate structure after etching.

\begin{figure}[H]
	\centering
	\begin{tikzpicture}[node distance = 3cm, auto, thick,scale=\CrossSectionOnly, every node/.style={transform shape}]
		\input{tikz_process_steps/fox.a.tex}
\fill[gateoxide] (0,2.75) rectangle (0.5,3.05);
\filldraw[line width=0, gateoxide] (0.5,2.75)--(1.25,2.0)--(1.25,2.3)--(0.5,3.05);
\fill[gateoxide] (1.25,2.0) rectangle (8.25,2.3);
\filldraw[line width=0, gateoxide] (9.00,2.75)--(8.25,2.0)--(8.25,2.3)--(9.00,3.05);
\fill[gateoxide] (9.00,2.75) rectangle (11.0,3.05);
\filldraw[line width=0, gateoxide] (11.0,2.75)--(11.75,2.0)--(11.75,2.3)--(11.0,3.05);
\fill[gateoxide] (11.75,2.0) rectangle (18.75,2.3);
\filldraw[line width=0, gateoxide] (19.5,2.75)--(18.75,2.0)--(18.75,2.3)--(19.5,3.05);
\fill[gateoxide] (19.5,2.75) rectangle (20.0,3.05);
	\end{tikzpicture} \\
	\includegraphics[scale=0.01]{down_arrow.png} \\
	\begin{tikzpicture}[node distance = 3cm, auto, thick,scale=\CrossSectionOnly, every node/.style={transform shape}]
		\coveringlayer{poly}{1.40}{0.00}
\fill[poly] (21.90,\SONOStopTHREE) rectangle (22.90,\polytop+0.2);

\input{tikz_process_steps/fox.a.tex}
\fill[gateoxide] (0,2.75) rectangle (0.5,3.05);
\filldraw[line width=0, gateoxide] (0.5,2.75)--(1.25,2.0)--(1.25,2.3)--(0.5,3.05);
\fill[gateoxide] (1.25,2.0) rectangle (8.25,2.3);
\filldraw[line width=0, gateoxide] (9.00,2.75)--(8.25,2.0)--(8.25,2.3)--(9.00,3.05);
\fill[gateoxide] (9.00,2.75) rectangle (11.0,3.05);
\filldraw[line width=0, gateoxide] (11.0,2.75)--(11.75,2.0)--(11.75,2.3)--(11.0,3.05);
\fill[gateoxide] (11.75,2.0) rectangle (18.75,2.3);
\filldraw[line width=0, gateoxide] (19.5,2.75)--(18.75,2.0)--(18.75,2.3)--(19.5,3.05);
\fill[gateoxide] (19.5,2.75) rectangle (20.0,3.05);

	\end{tikzpicture}
	\caption{Polysilicon}
\end{figure}

We use the LPCVD machine and deposit a layer of around 600nm polysilicon\footnote{\url{https://people.rit.edu/lffeee/LPCVD_Recipes.pdf}}.

We set the temperatue to 650\degreesC, the gas will be Silane ($Si H_4$ ($Si + 2H_2$)), the pressure will be set to 300 mTorr with a flow of 90sccm.

This will give us a growth rate of roughly 23.5 nm per minute, so for 600nm we let it grow half an hour.

\subsection{Patterning}

The resist is being deposited using spin coating and then baked depending on the baking time for the specific resist.
The layout for being exposed onto the resist is being extracted from the "poly" layer within the GDS2 file onto a bright field mask.

\begin{figure}[H]
	\centering
	\begin{tikzpicture}[node distance = 3cm, auto, thick,scale=\CrossSectionOnly, every node/.style={transform shape}]
		\coveringlayer{poly}{1.40}{0.00}
\fill[poly] (21.90,\SONOStopTHREE) rectangle (22.90,\polytop+0.2);

\input{tikz_process_steps/gate.gate_oxid_depositon.b.tex}

	\end{tikzpicture} \\
	\includegraphics[scale=0.01]{down_arrow.png} \\
	\begin{tikzpicture}[node distance = 3cm, auto, thick,scale=\CrossSectionOnly, every node/.style={transform shape}]
		\input{tikz_process_steps/gate.patterning.b.tex}
	\end{tikzpicture}
	\caption{Resist}
\end{figure}

\subsection{Etching}

\begin{figure}[H]
	\centering
	\begin{tikzpicture}[node distance = 3cm, auto, thick,scale=\CrossSectionOnly, every node/.style={transform shape}]
		\fill[resist,opacity=0.5] (5,3) rectangle (15.0,5.0);
\fill[resist] (5,3) rectangle (6.5,5.0);
\fill[resist] (13.5,3) rectangle (15.0,5.0);

\coveringlayer{poly}{1.40}{0.00}
\fill[poly] (21.90,\SONOStopTHREE) rectangle (22.90,\polytop+0.2);

\input{tikz_process_steps/gate.gate_oxid_depositon.b.tex}

	\end{tikzpicture} \\
	\includegraphics[scale=0.01]{down_arrow.png} \\
	\begin{tikzpicture}[node distance = 3cm, auto, thick,scale=\CrossSectionOnly, every node/.style={transform shape}]
		\input{tikz_process_steps/gate.etching.b.tex}
	\end{tikzpicture}
	\caption{Resist}
\end{figure}

\subsection{Cleaning}

\begin{figure}[H]
	\centering
	\begin{tikzpicture}[node distance = 3cm, auto, thick,scale=\CrossSectionOnly, every node/.style={transform shape}]
		\input{tikz_process_steps/gate.etching.b.tex}
	\end{tikzpicture} \\
	\includegraphics[scale=0.01]{down_arrow.png} \\
	\begin{tikzpicture}[node distance = 3cm, auto, thick,scale=\CrossSectionOnly, every node/.style={transform shape}]
		\def\welldepthAfox{0.75}
\def\welldepthBfox{1.25}
\def\welldepthCfox{1.5}

\newcommand{\stopper}[1]{
\filldraw[line width=0, isolationoxide] (#1,2.0) -- (#1+0.75,2.0) -- (#1+0.75,2.75);
\fill[isolationoxide] (#1+0.75,2.0) rectangle (#1+1.25,2.75);
\filldraw[line width=0, isolationoxide] (#1+1.25,2.75) -- (#1+1.25,2.0) -- (#1+2.0,2.0);
}

\newcommand{\bjtstopper}[1]{
\filldraw[line width=0, isolationoxide] (#1,2.0) -- (#1+0.25,2.0) -- (#1+0.25,2.75);
\fill[isolationoxide] (#1+0.25,2.0) rectangle (#1+0.5,2.75);
\filldraw[line width=0, isolationoxide] (#1+0.5,2.75) -- (#1+0.5,2.0) -- (#1+0.75,2.0);
}

% oxide
\fill[isolationoxide] (0,1.25) rectangle (55.0,2.0);

\fill[isolationoxide] (0,2.0) rectangle (0.75,2.75);
\filldraw[line width=0, isolationoxide] (0.75,2.75) -- (0.75,2.0) -- (1.5,2.0);

%separates n+ from p+
\stopper{2.5}
\stopper{8.0}

%separates n+ from p+
\stopper{13.5}

\stopper{16.5}
\stopper{17.0}

\stopper{20.0}

\stopper{24.5}
\stopper{25.0}

\bjtstopper{27.75}
\bjtstopper{29.0}
\bjtstopper{30.5}
\bjtstopper{31.75}

\stopper{33.5}
\stopper{33.75}

\bjtstopper{36.5}
\bjtstopper{37.75}
\bjtstopper{39.0}
\bjtstopper{40.25}

\filldraw[line width=0, isolationoxide] (41.75,2.0) -- (42.5,2.0) -- (42.5,2.75);
\fill[isolationoxide] (42.5,2.0) rectangle (55.0,2.75);

\input{tikz_process_steps/sti.a.tex}

% substrate
\fill[substrate] (0,0) rectangle (55,1.25);
\node at (2,0.5) {Silicon substrate};

% normal wells
\shade[upper left = nwell, upper right = nwell, lower right = substrate, lower left = substrate,] (1.25,\welldepthAfox) rectangle (8.25,2.0);
\shade[upper left = pwell, upper right = pwell, lower right = substrate, lower left = substrate,] (9.75,\welldepthAfox) rectangle (16.75,2.0);
\shade[upper left = nwell, upper right = nwell, lower right = substrate, lower left = substrate,] (18.25,\welldepthAfox) rectangle (25.25,2.0);
\shade[upper left = nwell, upper right = nwell, lower right = substrate, lower left = substrate,] (26.75,\welldepthAfox) rectangle (33.75,2.0);
\shade[upper left = nwell, upper right = nwell, lower right = substrate, lower left = substrate,] (35.25,\welldepthAfox) rectangle (42.25,2.0);

% p base
\shade[upper left = pbase, upper right = pbase, lower right = nwell, lower left = nwell,] (18.5,\welldepthBfox) rectangle (25.0,2.0);

% npn - pbase
\shade[upper left = pbase, upper right = pbase, lower right = nwell, lower left = nwell,] (28.25,\welldepthBfox) rectangle (32.0,2.0);

% pnp emitter and collector ring
\shade[upper left = pbase, upper right = pbase, lower right = nwell, lower left = nwell,] (35.5,\welldepthBfox) rectangle (37.0,2.0);
\shade[upper left = pbase, upper right = pbase, lower right = nwell, lower left = nwell,] (38.0,\welldepthBfox) rectangle (39.5,2.0);
\shade[upper left = pbase, upper right = pbase, lower right = nwell, lower left = nwell,] (40.5,\welldepthBfox) rectangle (42.0,2.0);

% n base for SONOS
\shade[upper left = nbase, upper right = nbase, lower right = pbase, lower left = pbase,] (18.75,\welldepthCfox) rectangle (24.75,2.0);

% n base for NPN emitter
\shade[upper left = nbase, upper right = nbase, lower right = pbase, lower left = pbase,] (29.5,\welldepthCfox) rectangle (30.75,2.0);


\fill[gateoxide] (5.0,2) rectangle (6.5,2.3);
\fill[gateoxide] (13.5,2) rectangle (15,2.3);
\fill[poly] (5.0,2.3) rectangle (6.5,3);
\fill[poly] (13.5,2.3) rectangle (15,3);

\fill[gateoxide,opacity=0.5] (5.0,2.0) rectangle (8.25,2.3);
\filldraw[line width=0, gateoxide,opacity=0.5] (9.00,2.75)--(8.25,2.0)--(8.25,2.3)--(9.00,3.05);
\fill[gateoxide,opacity=0.5] (9.00,2.75) rectangle (11.0,3.05);
\filldraw[line width=0, gateoxide,opacity=0.5] (11.0,2.75)--(11.75,2.0)--(11.75,2.3)--(11.0,3.05);
\fill[gateoxide,opacity=0.5] (11.75,2.0) rectangle (15.00,2.3);

\fill[poly,opacity=0.5] (5.0,2.3) rectangle (8.25,3.0);
\filldraw[line width=0, poly,opacity=0.5] (9.00,3.05)--(8.25,2.3)--(8.25,3.0)--(9.00,3.75);
\fill[poly,opacity=0.5] (9.00,3.05) rectangle (11.0,3.75);
\filldraw[line width=0, poly,opacity=0.5] (11.0,3.05)--(11.75,2.3)--(11.75,3.0)--(11.0,3.75);
\fill[poly,opacity=0.5] (11.75,2.3) rectangle (15.0,3.0);
	\end{tikzpicture}
	\caption{Resist}
\end{figure}
\newpage
\section{First vias (Contacts)}\label{via1}

This first set of vias connect the first metal layer to the active area.
These vias are in the fringe between front-end and back-end process.

\begin{figure}[H]
	\centering
	\begin{tikzpicture}[node distance = 3cm, auto, thick,scale=\CrossSectionOnly, every node/.style={transform shape}]
		\fill[isolationoxide] (0,\LowerMetal) rectangle (55,\LowerMoreMetal);

\fill[nitride] (0.0,\LowerMetal) rectangle (55.0,\LowerMetal+0.5);
\fill[isolationoxide] (0.0,2.0) rectangle (55.0,\LowerMetal+0.25);

\paintscaledmetal{nitride}{0.25}{0.5}
\paintscaledmetal{isolationoxide}{0.0}{0.25}

\input{tikz_process_steps/silicification.a.tex}

\paintcontacts{brown}{gray}{brown}


\paintscaledvias{white}{\UpperMetal}{\LowerMoreMetal}{0}


	\end{tikzpicture}
	\caption{Contact geometry target}
	\label{via1_cross_sections}
\end{figure}

As can be seen in \autoref{via1_cross_sections} the goal of this step is purely to deposit a layer of isolation oxide, get the holes into it, down to the silicide and polyside in order to form wires later on.
After this step the whole surface will still be covered by Aluminum, that's why the top view has been omitted.for this step.

\begin{figure}[H]
	\centering
	\begin{tikzpicture}[node distance =1cm, auto, thick,scale=\VLSILayout, every node/.style={transform shape}]
		\input{tikz_process_steps/pimplant.layout.tex}
% gate metal
\fill[gatemetal,opacity=\OpacityLayout] (4.8,1.75) rectangle (6.7,8);
\fill[gatemetal,opacity=\OpacityLayout] (13.3,1.75) rectangle (15.2,8);
\fill[gatemetal,opacity=\OpacityLayout] (4.8,8) rectangle (15.2,10);

%vias
\fill[via1,opacity=\OpacityLayout] (7,3.75) rectangle (7.5,4.25);
\fill[via1,opacity=\OpacityLayout] (12.5,3.75) rectangle (13,4.25);
\fill[via1,opacity=\OpacityLayout] (9.75,8.75) rectangle (10.25,9.25);
\fill[via1,opacity=\OpacityLayout] (2.25,3.75) rectangle (2.75,4.25);
\fill[via1,opacity=\OpacityLayout] (3.75,3.75) rectangle (4.25,4.25);
\fill[via1,opacity=\OpacityLayout] (15.75,3.75) rectangle (16.25,4.25);
\fill[via1,opacity=\OpacityLayout] (17.25,3.75) rectangle (17.75,4.25);

	\end{tikzpicture}
	\caption{First via layout}
	\label{via1_layout}
\end{figure}

It should be noted again that the via placement and dimensions in \autoref{via1_layout} are solely for demonstration purposes for the process only and are in not way the actual standard cell design for the final standard cell lib. \\

In a later iteration of this process we might be switching to Tungsten as the metal for this step.

\newpage

\subsection{Isolation dioxide layer}

We now need to grow a layer of thick oxide in order to isolate the Aluminum interconnect layer from the active area.

\begin{figure}[H]
	\centering
	\begin{tikzpicture}[node distance = 3cm, auto, thick,scale=\CrossSectionOnly, every node/.style={transform shape}]
		\fill[nitride] (0.0,\LowerMetal) rectangle (55.0,\LowerMetal+0.5);
\fill[isolationoxide] (0.0,2.0) rectangle (55.0,\LowerMetal+0.25);

\paintscaledmetal{nitride}{0.25}{0.5}
\paintscaledmetal{isolationoxide}{0.0}{0.25}

\input{tikz_process_steps/silicification.a.tex}

\paintcontacts{brown}{gray}{brown}


	\end{tikzpicture}
	\drawStepArrow{LTO deposition}
	\begin{tikzpicture}[node distance = 3cm, auto, thick,scale=\CrossSectionOnly, every node/.style={transform shape}]
		\fill[isolationoxide] (0,0) rectangle (20,\LowerMetal1);
\input{tikz_process_steps/pimplant.a.tex}

\filldraw[line width=0, isolationoxide] (5.75,2.0) -- (5.5,2.0) -- (5.75,3.0);
\filldraw[line width=0, isolationoxide] (6.75,2.0) -- (7.0,2.0) -- (6.75,3.0);

\filldraw[line width=0, isolationoxide] (11.25,2.0) -- (11.0,2.0) -- (11.25,3.0);
\filldraw[line width=0, isolationoxide] (12.5,2.0) -- (12.25,2.0) -- (12.25,3.0);

\filldraw[line width=0, isolationoxide] (22.75,2.4) -- (22.6,2.4) -- (22.75,3.4);
\filldraw[line width=0, isolationoxide] (23.75,2.4) -- (23.9,2.4) -- (23.75,3.4);

\fill[silicide] (1.5,1.9) rectangle (2.5,2);
\fill[silicide] (4.5,1.9) rectangle (5.5,2);
\fill[silicide] (5.75,2.9) rectangle (6.75,3.0);
\fill[silicide] (7,1.9) rectangle (8,2);

\fill[silicide] (10.0,1.9) rectangle (11.0,2);
\fill[silicide] (11.25,2.9) rectangle (12.25,3.0);
\fill[silicide] (12.5,1.9) rectangle (13.5,2);
\fill[silicide] (15.5,1.9) rectangle (16.5,2);

\fill[silicide] (19.0,1.9) rectangle (20.0,2);

\fill[silicide] (22.0,1.9) rectangle (22.6,2.0);
\fill[silicide] (22.75,3.3) rectangle (23.75,3.4);
\fill[silicide] (23.9,1.9) rectangle (24.5,2.0);

\fill[silicide] (27.0,1.9) rectangle (27.75,2.0);
\fill[silicide] (28.5,1.9) rectangle (29.0,2.0);
\fill[silicide] (29.75,1.9) rectangle (30.5,2.0);
\fill[silicide] (31.25,1.9) rectangle (31.75,2.0);
\fill[silicide] (32.5,1.9) rectangle (33.5,2.0);

\fill[silicide] (35.75,1.9) rectangle (36.5,2.0);
\fill[silicide] (37.25,1.9) rectangle (37.75,2.0);
\fill[silicide] (38.5,1.9) rectangle (39.0,2.0);
\fill[silicide] (39.75,1.9) rectangle (40.25,2.0);
\fill[silicide] (41.0,1.9) rectangle (41.75,2.0);

% diode contacts
\fill[silicide] (43.0,3.4) rectangle (44.0,3.5);
\fill[silicide] (47.0,3.4) rectangle (48.0,3.5);

% resistor contacts
\fill[silicide] (49.0,3.4) rectangle (49.5,3.5);
\fill[silicide] (53.5,3.4) rectangle (54.0,3.5);

	\end{tikzpicture}
	\caption{Oxide layer}
\end{figure}

\subsection{Pattering}

The resist is being deposited using spin coating and then baked depending on the baking time for the specific resist.
The layout for being exposed onto the resist is being extracted from the "n\_plus\_select" layer within the GDS2 file onto a \textbf{bright field} mask.
The requirement is a \textbf{negative} tone resist.

\begin{figure}[H]
	\centering
	\begin{tikzpicture}[node distance = 3cm, auto, thick,scale=\CrossAndTopSection, every node/.style={transform shape}]
		\input{tikz_process_steps/via1.patterning.a.tex}
	\end{tikzpicture}
	\begin{tikzpicture}[node distance = 3cm, auto, thick,scale=\CrossAndTopSection, every node/.style={transform shape}]
		\input{tikz_process_steps/via1.patterning.at.tex}
	\end{tikzpicture}
	\drawStepArrow{Mask: nselect}
	\begin{tikzpicture}[node distance = 3cm, auto, thick,scale=\CrossAndTopSection, every node/.style={transform shape}]
		\input{tikz_process_steps/via1.patterning.b.tex}
	\end{tikzpicture}
	\begin{tikzpicture}[node distance = 3cm, auto, thick,scale=\CrossAndTopSection, every node/.style={transform shape}]
		\input{tikz_process_steps/via1.patterning.bt.tex}
	\end{tikzpicture}
	\caption{N+ region resist mask}
\end{figure}

The thickness of the resist layer and the baking duration will variate depending on the specific equipment for which this process will be implemented with.
Also after the exposure and development, the hard baking shouldn't be forgotten!

\subsection{Etching}

We now need to open a window in the dioxide layer, through which we will inject carrier atoms into the silicon crystal structure.

\begin{figure}[H]
	\centering
	\begin{tikzpicture}[node distance = 3cm, auto, thick,scale=\CrossAndTopSection, every node/.style={transform shape}]
		\input{tikz_process_steps/via1.etching.a.tex}
	\end{tikzpicture}
	\begin{tikzpicture}[node distance = 3cm, auto, thick,scale=\CrossAndTopSection, every node/.style={transform shape}]
		\input{tikz_process_steps/via1.etching.at.tex}
	\end{tikzpicture}
	\drawStepArrow{}
	\begin{tikzpicture}[node distance = 3cm, auto, thick,scale=\CrossAndTopSection, every node/.style={transform shape}]
		\input{tikz_process_steps/via1.etching.b.tex}
	\end{tikzpicture}
	\begin{tikzpicture}[node distance = 3cm, auto, thick,scale=\CrossAndTopSection, every node/.style={transform shape}]
		\input{tikz_process_steps/via1.etching.bt.tex}
	\end{tikzpicture}
	\caption{N+ region opened}
\end{figure}

Since the silicon dioxide layer is 100nm thick and we wanna reach the silicon below we can use wet etching as described in the chemistry chapter.\\

\textbf{Possible approaches}:
\begin{itemize}
	\item \textbf{"AOE Etcher (DRY-AOE)" from HKUST} \\
	We can use anisotropic plasma etching for sharper borders.
	\item \textbf{Chemical solution} \\
	We can use buffered hydrofluoric acid (BOE (1:6)) at room temperature for a little bit over 3 minutes in order to get through the 100nm of oxide.\\
	Too long over 1 minutes might cause under-etch however!
\end{itemize}

\subsection{Cleaning}
\begin{figure}[H]
	\centering
	\begin{tikzpicture}[node distance = 3cm, auto, thick,scale=\CrossSectionOnly, every node/.style={transform shape}]
		\input{tikz_process_steps/via1.cleaning.a.tex}
	\end{tikzpicture}
	\drawStepArrow{}
	\begin{tikzpicture}[node distance = 3cm, auto, thick,scale=\CrossSectionOnly, every node/.style={transform shape}]
		\input{tikz_process_steps/via1.cleaning.b.tex}
	\end{tikzpicture}
	\caption{Resist removal}
\end{figure}
\newpage
\subsection{First metal layer}\label{metal1}
\begin{figure}[H]
	\centering
	\begin{tikzpicture}[node distance =1cm, auto, thick,scale=\VLSILayout, every node/.style={transform shape}]
		\input{tikz_process_steps/gate.layout.tex}
%vias
\fill[via1,opacity=\OpacityLayout] (7,3.75) rectangle (7.5,4.25);
\fill[via1,opacity=\OpacityLayout] (12.5,3.75) rectangle (13,4.25);
\fill[via1,opacity=\OpacityLayout] (9.75,8.75) rectangle (10.25,9.25);
\fill[via1,opacity=\OpacityLayout] (2.25,3.75) rectangle (2.75,4.25);
\fill[via1,opacity=\OpacityLayout] (3.75,3.75) rectangle (4.25,4.25);
\fill[via1,opacity=\OpacityLayout] (15.75,3.75) rectangle (16.25,4.25);
\fill[via1,opacity=\OpacityLayout] (17.25,3.75) rectangle (17.75,4.25);

\fill[Fuchsia,opacity=0.2] (9,8) rectangle (11,12);
\fill[Fuchsia,opacity=0.2] (6.5,3) rectangle (13.5,5);
\fill[Fuchsia,opacity=0.2] (9,0) rectangle (11,3);
\fill[Fuchsia,opacity=0.2] (2,0) rectangle (4.5,12);
\fill[Fuchsia,opacity=0.2] (15.5,0) rectangle (18,12);
	\end{tikzpicture}
	\caption{First metal layout}
	\label{metal1_layout}
\end{figure}
\newpage
\subsection{Additional vias}\label{via2}
\begin{figure}[H]
	\centering
	\begin{tikzpicture}[node distance = 3cm, auto, thick,scale=\VLSILayout, every node/.style={transform shape}]
		\input{tikz_process_steps/via1.layout.tex}
\fill[Fuchsia,opacity=0.2] (9,8) rectangle (11,12);
\fill[Fuchsia,opacity=0.2] (6.5,3) rectangle (13.5,5);
\fill[Fuchsia,opacity=0.2] (9,0) rectangle (11,3);
\fill[Fuchsia,opacity=0.2] (2,0) rectangle (4.5,12);
\fill[Fuchsia,opacity=0.2] (15.5,0) rectangle (18,12);
	\end{tikzpicture}
	\caption{Additional via layout}
	\label{via2_layout}
\end{figure}
\newpage
\subsection{Additional metal layer}\label{metal2}
\begin{figure}[H]
	\centering
	\begin{tikzpicture}[node distance = 3cm, auto, thick,scale=\VLSILayout, every node/.style={transform shape}]
		\input{tikz_process_steps/metal1.layout.tex}
	\end{tikzpicture}
	\caption{Additional metal layout}
	\label{metal2_layout}
\end{figure}

\newpage
\section{Testing}
In order to get an idea on how strongly the actual product differs from the mathematical models being used before hand to define the initial parameters, one has to run a bunch of test wafers through the process over and over again, tweak parameters and measure out all kinds of aspects of the device.

This is also important for gathering the data which will finally end up within the data sheet.

In this chapter we will tackle the multiple different measurement and test objects we will need to put into the seal area during production in order to ensure the consistent quality of the chips we're going to sell in the end.\\

\textbf{Testing parameters:}
\begin{itemize}
	\item The capacity per $m^2$ from the n-well to the gate electrode
	\item The capacity per $m^2$ from  the p-well to the gate electrode
	\item The actual resistance/m of the p-well
	\item The actual resistance/m of the n-well
	\item The actual resistance/m of the poly layer
	\item The actual resistance/m of the silicide+p-implant layer
	\item The actual resistance/m of the silicide+n-implant layer
	\item The actual resistance/m of the silicide+poly layer
	\item NMOS as well as PMOS:
	\begin{itemize}
		\item Frequency characteristics (Transient curve)
		\item Drain-source resistance vs. gate-source voltage
		\item Actual threshold voltage
	\end{itemize}
\end{itemize}

Each of these values will variate based on the location on the wafer, because of diffraction towards the edge and also because of the uneven nature of the wafer.

Measuring the same test structure multiple times placed at multiple locations on the wafer will allow us to calculate an effective average value and tolerance ranges for taking into account for further dimensioning of circuits..

\newpage
\section{Design rules}
\section{Layer Definitions}\label{design_rules_layer_definitions}

\begin{center}
%    \captionof{table}{Mask Layers (for Production)}
    \begin{tabular}{|l|c|c|p{7.5cm}|}
    \hline
        \textbf{Name} & \textbf{GDSII} & \textbf{CIF} & \textbf{Description} \\
    \hline
        PWELL   & 41 & CWP & p-well \\
        NWELL   & 42 & CWN & n-well \\
        ACTIVE  & 43 & CAA & active area \\
        PPLUS   & 44 & CSP & $p^{+}$ implant \\
        NPLUS   & 45 & CSN & $n^{+}$ implant \\
        POLY    & 46 & CPG & poly silicium \\
        CONTACT & 25 & CCC & contact (connects METAL1 to POLY) \\
        METAL1  & 49 & CM1 & lowest metal layer \\
        VIA1    & 50 & CV1 & via layer (connects METAL2 to METAL1) \\
        METAL2  & 51 & CM2 & second metal layer \\
        VIA2    & 61 & CV2 & via layer (connects METAL3 to METAL2) \\
        METAL3  & 62 & CM3 & third metal layer \\
        GLASS   & 52 & COG & passivation / isolation \\
    \hline
    \end{tabular}
\end{center}

\section{General Requirements}\label{design_rules_general_requirements}

\begin{flushleft}
    \begin{tabular}{c p{9cm} c}
        1.1  & All scaled dimensions are specified in Lambda & $\lambda$. \\
        1.2  & All fixed dimensions are specified in Microns & $\mu m$. \\
        1.3  & All geometries must be drawn on grid. The grid size is & $1 \lambda$. \\
        1.4  & Polygones should be rectangles with 90 degree angles only. \\
        1.5  & The die size should be an integer multiple of & $10 \mu m$. \\
    \end{tabular}
\end{flushleft}

\section{Process Layer Overview}\label{design_rules_process_layer_overview}

\begin{center}
%    \captionof{table}{Mask Layer Minimum Values}
    \begin{tabular}{|l|c|c|}
    \hline
        \textbf{Name} & \textbf{Minimum Width} & \textbf{Minimum Spacing} \\
    \hline
        PWELL   & 10 $\lambda$ & 10 $\lambda$ \\
        NWELL   & 10 $\lambda$ & 10 $\lambda$ \\
        ACTIVE  & 3  $\lambda$ & 3  $\lambda$ \\
        POLY    & 2  $\lambda$ & 2  $\lambda$ \\
        CONTACT & 2  $\lambda$ & 2  $\lambda$ \\
        METAL1  & 4  $\lambda$ & 4  $\lambda$ \\
        VIA1    & 2  $\lambda$ & 3  $\lambda$ \\
        METAL2  & 4  $\lambda$ & 4  $\lambda$ \\
        VIA2    & 2  $\lambda$ & 3  $\lambda$ \\
        METAL3  & 6  $\lambda$ & 4  $\lambda$ \\
    \hline
    \end{tabular}
\end{center}

\section{Structure Rules}\label{design_rules_structure_rules}

\subsection{PWELL Rules}\label{design_rules_pwell_rules}

\begin{center}
    \begin{tikzpicture}[scale=(.33)]
        \filldraw[fill=pwell] (0, 0) rectangle ++(10, 10) node[near end] {PWELL};
        \draw[<->] (0, 1) -- (10, 1) node[near start, above] {2.2.1};
        \draw[<->] (1, 0) -- (1, 10) node[near start, left] {2.2.1};
        \draw[<->] (10, 5) -- (19, 5) node[near start, above] {2.2.2};
        \filldraw[fill=pwell] (19, 0) rectangle ++(10, 10) node[near end] {PWELL};
        \draw[<->] (29, 5) -- (35, 5) node[near start, above] {2.2.3};
        \filldraw[fill=pwell] (35, 0) rectangle ++(10, 10) node[near end] {PWELL};
    \end{tikzpicture}
\end{center}

\begin{flushleft}
    \begin{tabular}{c l c}
        2.1.1  & Minimum Width of PWELL is & $ 10 \lambda $. \\
        2.1.2  & Minimum Spacing to PWELL at different potential is & $ 10 \lambda $. \\
        2.1.3  & Minimum Spacing to PWELL at same potential is & $ 0 \lambda $. \\
        2.1.4  & Minimum Spacing to NWELL is & $ 12 \lambda $. \\
    \end{tabular}
\end{flushleft}

\subsection{NWELL Rules}\label{design_rules_nwell_rules}

\begin{center}
    \begin{tikzpicture}[scale=(.33)]
        \filldraw[fill=nwell] (0, 0) rectangle ++(10, 10) node[near end] {NWELL};
        \draw[<->] (0, 1) -- (10, 1) node[near start, above] {2.2.1};
        \draw[<->] (1, 0) -- (1, 10) node[near start, left] {2.2.1};
        \draw[<->] (10, 5) -- (19, 5) node[near start, above] {2.2.2};
        \filldraw[fill=nwell] (19, 0) rectangle ++(10, 10) node[near end] {NWELL};
        \draw[<->] (29, 5) -- (35, 5) node[near start, above] {2.2.3};
        \filldraw[fill=nwell] (35, 0) rectangle ++(10, 10) node[near end] {NWELL};
    \end{tikzpicture}
\end{center}

\begin{flushleft}
    \begin{tabular}{c l c}
        2.2.1  & Minimum Width of NWELL is & $ 10 \lambda $. \\
        2.2.2  & Minimum Spacing to NWELL at different potential is & $ 10 \lambda $. \\
        2.2.3  & Minimum Spacing to NWELL at same potential is & $ 0 \lambda $. \\
        2.2.4  & Minimum Spacing to PWELL is & $ 12 \lambda $. \\
    \end{tabular}
\end{flushleft}

\subsection{ACTIVE Rules}\label{design_rules_active_rules}

\begin{flushleft}
    \begin{tabular}{c l c}
        2.3.1   & Minimum Width of ACTIVE is & $3 \lambda$. \\
        2.3.2   & Minimum Spacing to ACTIVE is & $3 \lambda$. \\
        2.3.3   & Minimum Source/Drain surround by PWELL/NWELL is & $6 \lambda$. \\
        2.3.4   & Minimum Substrate/NWELL contact sourround by PWELL/NWELL is & $3 \lambda$. \\
        2.3.5   & Minimum Spacing to ACTIVE of opposite type is & $4 \lambda$. \\
    \end{tabular}
\end{flushleft}

\subsection{POLY Rules}\label{design_rules_poly_rules}

\begin{center}
    \begin{tikzpicture}[scale=(.4)]
        \filldraw[fill=poly] (0, 0) rectangle ++(2, 12) node[near end] {POLY};
        \draw[<->] (0, 6) -- (2, 6) node[near start, above] {2.4.1};
        \draw[<->] (2, 4) -- (4, 4) node[near start, above] {2.4.3};
        \filldraw[fill=nitride] (7, 0) rectangle ++(5, 12) node[near end] {ACTIVE};
        \draw[<->] (10, 5) -- (10, 7) node[near start, left] {2.4.2};
        \draw[<->] (9, 0) -- (9, 3) node[near start, left] {2.4.5};
        \filldraw[fill=poly] (4, 5) -- ++(10, 0) -- ++(0, -2) -- ++(-8, 0) -- ++(0, -3) -- ++(-2, 0) -- ++(0, 5) node[near end] {POLY};
        \filldraw[fill=poly] (4, 7) -- ++(0, 5) -- ++(2, 0) -- ++(0, -3) -- ++(8, 0) -- ++(0, -2) -- ++(-10, 0) node[near end] {POLY};
        \draw[<->] (12, 8) -- (14, 8) node[near start, above] {2.4.4};
        \draw[<->] (6, 10) -- (7, 10) node[near start, above] {2.4.6};
    \end{tikzpicture}
\end{center}

\begin{flushleft}
    \begin{tabular}{c l c}
        2.4.1   & Minimum Width of POLY is & $2 \lambda$. \\
        2.4.2   & Minimum Spacing to POLY over ACTIVE is & $2 \lambda$. \\
        2.4.3   & Minimum Spacing to POLY over others is & $2 \lambda$. \\
        2.4.4   & Minimum Gate extension beyond ACTIVE is & $2 \lambda$. \\
        2.4.5   & Minimum ACTIVE extension beyond POLY is & $3 \lambda$. \\
        2.4.6   & Minimum Spacing of POLY to ACTIVE is & $1 \lambda$. \\
    \end{tabular}
\end{flushleft}

\subsection{NPLUS/PPLUS Rules}\label{design_rules_plus_rules}

\begin{flushleft}
    \begin{tabular}{c l c}
        2.5.1   & Minimum Spacing to POLY is & $3 \lambda$. \\
        2.5.2   & Minimum Overlap to ACTIVE is & $2 \lambda$. \\
        2.5.3   & Minimum Overlap to CONTACT is & $1 \lambda$. \\
    \end{tabular}
\end{flushleft}

\subsection{CONTACT Rules}\label{design_rules_contact_rules}

\begin{flushleft}
    \begin{tabular}{c l c}
        2.6.1   & Exact Width/Height of CONTACT is & $ 2 \lambda \times 2 \lambda$. \\
        2.6.2   & Minimum Overlap by POLY or ACTIVE is & $1 \lambda$. \\
        2.6.3   & Minimum Spacing to CONTACT is & $2 \lambda$. \\
        2.6.4   & Minimum Spacing to Gate is & $2 \lambda$. \\
        2.6.5   & Minimum Spacing of POLY CONTACT to other POLY is & $4 \lambda$. \\
        2.6.6   & Minimum Spacing of ACTIVE CONTACT to POLY CONTACT is & $4 \lambda$. \\
    \end{tabular}
\end{flushleft}

\subsection{METAL1 Rules}\label{design_rules_metal1_rules}

\begin{center}
    \begin{tikzpicture}[scale=(.33)]
        \filldraw[fill=metal1] (0, 0) rectangle ++(4, 10) node[near end] {METAL1};
        \draw[<->] (0, 1) -- (4, 1) node[near start, above] {2.7.1};
        \draw[<->] (4, 5) -- (8, 5) node[near start, above] {2.7.2};
        \filldraw[fill=metal1] (8, 0) rectangle ++(4, 10) node[near end] {METAL1};
        \filldraw[fill=titanium] (9, 1) rectangle ++(2, 2); % !! color code
        \draw[<->] (8, 2) -- (9, 2) node[near start, above] {2.7.3};
        \draw[<->] (12, 5) -- (18, 5) node[near start, above] {2.7.4};
        \filldraw[fill=metal1] (18, 0) rectangle ++(20, 10) node[near end] {(Big-)METAL1};
        \draw[<->] (18, 1) -- (38, 1) node[near start, above] {$ \geq 10 \mu m$};
    \end{tikzpicture}
\end{center}

\begin{flushleft}
    \begin{tabular}{c l c}
        2.7.1   & Minimum Width of METAL1 is & $4 \lambda$. \\
        2.7.2   & Minimum Spacing of METAL1 is & $4 \lambda$. \\
        2.7.3   & Minimum Overlap of CONTACT or VIA is & $1 \lambda$. \\
        2.7.4   & Minimum Spacing to METAL1 for lines wider than $10 \mu m$ is & $6 \lambda$. \\
    \end{tabular}
\end{flushleft}

\subsection{VIA1 Rules}\label{design_rules_via1_rules}

\begin{flushleft}
    \begin{tabular}{c l c}
        2.8.1   & Exact Width/Height of VIA1 is & $ 2 \lambda \times 2 \lambda$. \\
        2.8.2   & Minimum Spacing to VIA1 is & $3 \lambda$. \\
        2.8.3   & Minimum Spacing to CONTACT is & $2 \lambda$. \\
        2.8.4   & Minimum Spacing to POLY or ACTIVE is & $2 \lambda$. \\
    \end{tabular}
\end{flushleft}

\subsection{METAL2 Rules}\label{design_rules_metal2_rules}

\begin{center}
    \begin{tikzpicture}[scale=(.33)]
        \filldraw[fill=metal2] (0, 0) rectangle ++(4, 10) node[near end] {METAL2};
        \draw[<->] (0, 1) -- (4, 1) node[near start, above] {2.9.1};
        \draw[<->] (4, 5) -- (8, 5) node[near start, above] {2.9.2};
        \filldraw[fill=metal2] (8, 0) rectangle ++(4, 10) node[near end] {METAL2};
        \filldraw[fill=via1] (9, 1) rectangle ++(2, 2); % !! color code
        \draw[<->] (8, 2) -- (9, 2) node[near start, above] {2.9.3};
        \draw[<->] (12, 5) -- (18, 5) node[near start, above] {2.9.4};
        \filldraw[fill=metal2] (18, 0) rectangle ++(20, 10) node[near end] {(Big-)METAL2};
        \draw[<->] (18, 1) -- (38, 1) node[near start, above] {$ \geq 10 \mu m$};
    \end{tikzpicture}
\end{center}

\begin{flushleft}
    \begin{tabular}{c l c}
        2.9.1   & Minimum Width of METAL2 is & $4 \lambda$. \\
        2.9.2   & Minimum Spacing to METAL2 is & $4 \lambda$. \\
        2.9.3   & Minimum Overlap to VIA1 is & $1 \lambda$. \\
        2.9.4   & Minimum Spacing to METAL2 for lines wider than $10 \mu m$ is & $6 \lambda$. \\
    \end{tabular}
\end{flushleft}

\subsection{VIA2 Rules}\label{design_rules_via2_rules}

\begin{flushleft}
    \begin{tabular}{c l c}
	2.10.1   & Exact Width/Height of VIA2 is & $ 2 \lambda \times 2 \lambda$. \\
        2.10.2   & Minimum Spacing to VIA2 is & $3 \lambda$. \\
    \end{tabular}
\end{flushleft}

\subsection{METAL3 Rules}\label{design_rules_metal3_rules}

\begin{center}
    \begin{tikzpicture}[scale=(.33)]
        \filldraw[fill=metal2] (0, 0) rectangle ++(4, 10) node[near end] {METAL3}; % !! color code
        \draw[<->] (0, 1) -- (4, 1) node[near start, above] {2.11.1};
        \draw[<->] (4, 5) -- (8, 5) node[near start, above] {2.11.2};
        \filldraw[fill=metal2] (8, 0) rectangle ++(4, 10) node[near end] {METAL3}; % !! color code
        \filldraw[fill=via2] (9, 1) rectangle ++(2, 2); % !! color code
        \draw[<->] (8, 2) -- (9, 2) node[near start, above] {2.11.3};
        \draw[<->] (12, 5) -- (18, 5) node[near start, above] {2.11.4};
        \filldraw[fill=metal2] (18, 0) rectangle ++(20, 10) node[near end] {(Big-)METAL3}; % !! color code
        \draw[<->] (18, 1) -- (38, 1) node[near start, above] {$ \geq 10 \mu m$};
    \end{tikzpicture}
\end{center}

\begin{flushleft}
    \begin{tabular}{c l c}
        4.11.1  & Minimum Width of METAL3 is & $6 \lambda$. \\
        2.11.2  & Minimum Spacing to METAL3 is & $4 \lambda$. \\
        2.11.3  & Minimum Overlap to VIA2 is & $1 \lambda$. \\
        2.11.4  & Minimum Spacing to METAL3 for lines wider than $10 \mu m$ is & $6 \lambda$. \\
    \end{tabular}
\end{flushleft}

\subsection{GLASS Rules}\label{design_rules_glass_rules}

\begin{flushleft}
    \begin{tabular}{c l c}
	2.12.1  & Exact Width/Height of GLASS is & $ 120 \lambda \times 120 \lambda$. \\
        2.12.2  & Minimum METAL3 Overlap of GLASS opening is & $ 12 \lambda $. \\
        2.12.3  & Minimum Spacing of pad METAL3 to unrelated METAL is & $30 \mu m$. \\
        2.12.4  & Minimum Spacing of pad METAL3 to ACTIVE or POLY is & $15 \mu m$. \\
    \end{tabular}
\end{flushleft}




\end{document}
